\documentclass[11pt,a4paper]{article}
\usepackage[T1]{fontenc}
\usepackage[utf8]{inputenc}
\usepackage{lmodern,microtype}
\usepackage[margin=24mm]{geometry}
\usepackage{amsmath,amssymb,amsthm,mathtools}
\usepackage{booktabs,array,tabularx,longtable}
\usepackage{xcolor,enumitem,fancyhdr,listings}
\usepackage[unicode,hypertexnames=false,colorlinks=true,linkcolor=blue!40!black,urlcolor=blue!50!black,citecolor=blue!40!black]{hyperref}
\definecolor{ink}{HTML}{153047}
\definecolor{quiet}{HTML}{516371}
\newtheorem{theorem}{Theorem}[section]
\newtheorem{lemma}[theorem]{Lemma}
\newtheorem{proposition}[theorem]{Proposition}
\newtheorem{corollary}[theorem]{Corollary}
\theoremstyle{remark}
\newtheorem{remark}[theorem]{Remark}
\DeclareMathOperator{\ad}{ad}
\DeclareMathOperator{\diam}{diam}
\DeclareMathOperator{\rank}{rank}
\newcommand{\E}{\mathbb E}
\newcommand{\one}{\mathbf 1}
\numberwithin{equation}{section}
\setlength{\parskip}{4pt}
\setlength{\parindent}{0pt}
\setlength{\emergencystretch}{2em}
\setlist{topsep=4pt,itemsep=2pt}
\pagestyle{fancy}\fancyhf{}
\fancyhead[L]{\small\color{quiet}Caporossi--Hansen: global classes}
\fancyhead[R]{\small\color{quiet}24 November 2026}
\fancyfoot[C]{\small\thepage}
\setlength{\headheight}{14pt}
\lstset{language=Python,basicstyle=\ttfamily\scriptsize,
 breaklines=true,columns=fullflexible,keepspaces=true,
 showstringspaces=false,frame=single,rulecolor=\color{quiet!35},
 aboveskip=8pt,belowskip=8pt,tabsize=4}
\hypersetup{
 pdftitle={The Caporossi--Hansen Additive Inequality: All Trees, All Leafless Graphs, and a Smaller General Cutoff},
 pdfauthor={Keio and Shaitetsu},
 pdfsubject={Complete class-wide proofs and exact finite certificates; the unrestricted conjecture remains unresolved}}
\begin{document}
\begin{center}
{\small\color{quiet}VERIFIABLE MATHEMATICAL REASONING / RESEARCH CONTINUATION}\\[9pt]
{\LARGE\bfseries\color{ink}The Caporossi--Hansen\\[3pt] Additive Inequality}\\[8pt]
{\Large All trees, all leafless graphs,\\[3pt]
and a smaller general cutoff}\\[10pt]
{\large Keio and Shaitetsu}\\[4pt]
24 November 2026
\end{center}
\vspace{3pt}\hrule\vspace{9pt}

\textbf{Abstract.}
We continue the explicit $n\ge10^{10}$ argument for
$R(G)-\ad(G)\ge\sqrt{n-1}-2+2/n$. We prove this inequality for every
tree, with equality exactly for stars. This proof is independent of the
corpus weighted-distance theorem: it combines an edge potential, a
degree-sequence extremum, and a sharp weighted-tree distance bound with
a prescribed heavy vertex. Infinite orders are reduced to a compact
rational interval certificate; the remaining orders are covered by
exhaustive integer parameter certificates and a small family of explicitly
classified trees.
Using only the uniform specialization of the corpus weighted theorem,
we also prove the strict additive inequality for every connected graph
of minimum degree at least two, and lower the sufficient order for
arbitrary connected graphs to $3\,000\,000$.
A smallest remaining counterexample of order at least $44$ cannot have
a leaf adjacent to a vertex of degree two.
The unrestricted conjecture is not proved: the remaining class consists
of graphs with both leaves and cycles below the new cutoff.
All proof certificates are included as source code in this single
manuscript.

\section{The conclusions and their precise dependencies}

All graphs are finite, simple and undirected. Unless stated otherwise,
they are connected and have order $n\ge2$. Put
\[
 R(G)=\sum_{uv\in E(G)}\frac1{\sqrt{d_ud_v}},\qquad
 M(G)=\sum_{uv\in E(G)}\frac1{\max(d_u,d_v)},\qquad
 \mu(G)=\binom n2^{-1}\sum_{\{u,v\}}d_G(u,v).
\]
Here $M$ is a modified Randi\'c sum, not a mode of the distance distribution.
The target, proposed by Caporossi and Hansen \cite{cap,favaron}, is
\[
 R(G)-\mu(G)\ge\gamma_n,\qquad
 \gamma_n=\sqrt{n-1}-2+\frac2n.
 \tag{CH}\label{eq:CH}
\]
The star has $R=\sqrt{n-1}$ and $\mu=2-2/n$, so it attains this value.

\begin{center}\small
\begin{tabularx}{\linewidth}{@{}>{\raggedright\arraybackslash}p{.25\linewidth}
 X >{\raggedright\arraybackslash}p{.24\linewidth}@{}}
\toprule
Class or statement & Result in this manuscript & Dependency\\
\midrule
Every tree, every order & \textup{(CH)}, with equality only on stars &
Independent of the corpus weighted theorem\\
Every graph without leaves & Strict \textup{(CH)} at every order &
Uniform corpus input \textup{(U)}\\
Every connected graph of order at least $3\,000\,000$ &
\textup{(CH)}, with equality only on stars &
Uniform corpus input \textup{(U)}\\
Smallest counterexample of order at least $44$ &
Every leaf has a neighbor of degree at least three &
Elementary deletion argument\\
Graphs with leaves and cycles below $3\,000\,000$ &
Not settled in full & The remaining unrestricted case\\
\bottomrule
\end{tabularx}
\end{center}

The earlier corpus argument used the sufficient order $10^{27}$
\cite{vmrc227}; the preceding report in this conversation reduced it
to $10^{10}$. The new cutoff and both unrestricted-in-order graph
classes are proved below. The computations in the proof are exhaustive
over the specified parameter domains, not tests on a chosen graph sample.
External peer review and historical priority have not been established.

\subsection{The sole imported input}

The corpus manuscript [227] records the weighted-distance theorem
\[
 \sum_{u,v}p_up_vd_G(u,v)+2\sum_vp_v^2\le M(G)+1
 \quad\text{for every probability vector }p.
 \tag{W}\label{eq:W-new}
\]
Uniform mass gives the only specialization needed for our leafless
and general-order results:
\begin{equation}
 M(G)\ge\frac{(n-1)\mu(G)-(n-2)}n.
 \tag{U}\label{eq:U-new}
\end{equation}
Its proof is in the frozen corpus source \cite{vmrc227}.
All other mathematical estimates used here are supplied in this report.
In particular the all-tree theorem does not use \textup{(U)} or
\textup{(W)}.

\section{Elementary degree and distance estimates}

\begin{lemma}[Degree-ratio bridge]\label{lem:new-bridge}
For a connected graph of order $n\ge2$, let $r=\sqrt{\Delta(G)}$. Then
\[
 n-2R\le(r-1)(R-M),\qquad
 (r+1)R\ge n+(r-1)M.
\]
Consequently $R\ge\sqrt{n-1}$. If the minimum degree is at least two,
the first inequality also holds with $r=\sqrt{\Delta(G)/2}$.
\end{lemma}
\begin{proof}
For endpoint degrees $a\le b$,
\[
 (a^{-1/2}-b^{-1/2})^2
 =(\sqrt{b/a}-1)((ab)^{-1/2}-b^{-1}).
\]
Sum over the edges and use
$\sum_{uv\in E}(d_u^{-1}+d_v^{-1})=n$.
The degree ratio is at most $\Delta$, or at most $\Delta/2$
when every degree is at least two.
The edges at a maximum-degree vertex contribute exactly one to $M$,
so $M\ge1$. Since $r\le\sqrt{n-1}$, the second inequality gives
\[
 R\ge1+\frac{n-2}{r+1}
 \ge1+\frac{n-2}{\sqrt{n-1}+1}=\sqrt{n-1}.
\]
\end{proof}

\begin{lemma}[Broom distance bound]\label{lem:new-broom}
For a connected graph of order $n\ge2$, write $d=\Delta(G)$ and
$k=n-1-d$. Then
\[
 \mu(G)\le B(n,d):=
 2-\frac2n+\frac{k(k+1)(3n-2k-7)}{3n(n-1)}.
\]
The bound is attained by a star of degree $d$ with one arm extended
by $k$ additional vertices. In particular
$\mu(G)\le(n+1)/3$.
\end{lemma}
\begin{proof}
Choose a vertex of degree $d$ and extend all its incident edges to a
spanning tree. Start with this $(d+1)$-vertex star and add the remaining
vertices, in a rooted order, one at a time as leaves. When the new order
is $t$, the new leaf's sorted distances $b_1\le\cdots\le b_{t-1}$
satisfy $b_i\le i$: a path to a vertex at distance $j$ contains all
smaller positive distances. They also satisfy $b_i\le t-d+1$,
because a tree containing a degree-$d$ vertex has diameter at most
$t-d+1$. A path through that vertex leaves at least $d-2$ of its
neighbors off the path; a path avoiding it gives an even stronger count.
Thus the new transmission is at most
\[
 \sum_{i=1}^{t-1}\min\{i,t-d+1\}
 =\binom t2-\binom{d-1}{2}.
\]
The original star has Wiener index $d^2$. Sum the last bound from
$t=d+2$ to $n$, divide by $\binom n2$, and substitute $d=n-1-k$.
This is $B(n,d)$, and a single extended arm attains every summand.
The case $d=1$ is $K_2$. Dropping the degree constraint in the same
leaf-addition proof gives the path bound $(n+1)/3$.
\end{proof}

\subsection{The exact weighted-tree estimate}



\begin{lemma}[Weighted tree with a heavy vertex]\label{lem:ch-heavy-tree}
Let a tree have $h\ge1$ edges and positive integer vertex weights totaling
$n$. Suppose one vertex has weight at least $a$. Put $P=n-h+1$ and
$A=\max\{a,\lceil P/2\rceil\}$. Then
\[
 \mathbb E d_T(X,Y)
 \le\frac h2-\frac{h(h^2-1)}{6n^2}
       -\frac{2h}{n^2}\left(A-\frac P2\right)^2.
\]
Here independent samples use the normalized vertex weights. The bound is
attained by a path with internal weights one and endpoint weights $A,P-A$.
In particular the same upper bound remains valid if the final squared term
is replaced by $\max\{0,a-P/2\}^2$. For $h=0$ the expectation is zero.
\end{lemma}

\begin{proof}
For an edge let $s_e$ be the smaller of its two component weights. The
expected distance is $2n^{-2}\sum_e s_e(n-s_e)$. Fix an integer $t\ge1$.
The edges with $s_e\ge t$, if nonempty, form a subtree $S$ with, say,
$m$ edges. The usual two outer-component count gives $n\ge2t+m-1$.
There is also the bound $n\ge a+t+m-1$. Indeed, if the heavy vertex lies
in a leaf's outer component of $S$, that component has weight at least
$a$, another leaf's outer component has weight at least $t$, and the
remaining vertices/components contribute at least $m-1$. If the heavy
vertex lies elsewhere, the two outer components of weight at least $t$
and the excess weight at the heavy vertex give an at least as strong
bound. More explicitly, for $\ell$ leaves the two cases give respectively
$a+(\ell-1)t+(m+1-\ell)$ and
$\ell t+(m+1-\ell)+(a-1)$, each at least $a+t+m-1$.
Thus the ordered cut sizes satisfy
\[
 s_j\le\min\left\{
       \left\lfloor\frac{n-j+1}{2}\right\rfloor,
       n-a-j+1\right\},\qquad 1\le j\le h.
\]
These bounds are attained simultaneously by the stated path. One way to
check this is to count its cuts with $s_e\ge t$: the cut positions are the
consecutive integers $A,A+1,\ldots,A+h-1$, and the count is
\[
 \max\{0,\min(h,n-2t+1,n-A-t+1)\}.
\]
When $A>a$, the last term is redundant because the endpoint weights differ
by at most one. Since $s(n-s)$ increases for $0\le s\le n/2$, it remains
to sum the path's cut products. Their deviations from $n/2$ have squared
sum
\[
 \sum_{i=0}^{h-1}(A+i-n/2)^2
 =\frac{h(h^2-1)}{12}+h(A-P/2)^2.
\]
Substitution proves the result. Positivity of the other $h$ vertex weights
ensures $a\le n-h$, so the asserted endpoint weights are positive.
\end{proof}

\section{The additive inequality for every tree}
\label{sec:ch-all-trees}

For a finite tree $T$ of order $n\ge2$, write
\[
 R(T)=\sum_{uv\in E(T)}\frac1{\sqrt{d_ud_v}},\qquad
 \mu(T)=\binom n2^{-1}\sum_{\{u,v\}\subseteq V(T)}d_T(u,v),
 \qquad
 \gamma_n=\sqrt{n-1}-2+\frac2n.
\]

\begin{theorem}\label{thm:ch-all-trees}
Every finite tree of order $n\ge2$ satisfies
\[
                  R(T)-\mu(T)\ge\gamma_n.
\]
Equality holds if and only if $T\cong K_{1,n-1}$.
\end{theorem}

The proof is independent of the corpus theorem \textup{(W)}. It uses the
elementary heavy-vertex weighted-tree lemma
(Lemma~\ref{lem:ch-heavy-tree}), two degree estimates proved below, and
finite certificates whose comparisons are entirely rational or integral.
The infinite range is reduced to a fixed compact two-dimensional domain;
the finite-order calculation checks parameter bounds, with only seven
small parameter profiles requiring a direct enumeration of tree shapes.

The tree of order two is $K_2$, and direct evaluation gives equality.
For the following notation assume $n\ge3$. Throughout the proof, let
$L$ be the number of leaves, $I=n-L$ the number
of nonleaves, and $\Delta$ the maximum degree. Put
\[
 N=n-1,\qquad h=I-1=N-L,\qquad k=N-\Delta,
 \qquad s=\sqrt N,\qquad r=\sqrt\Delta.
\]
Every vertex of degree $d$ separates a tree into $d$ branches, each
containing a distinct leaf. Hence $\Delta\le L$ and $k\ge h$.
A tree with $h=0$ is a star. Its values are
$R=\sqrt N$ and $\mu=2-2/n$, so it attains equality.
We henceforth consider nonstars, for which $n\ge4$ and $h\ge1$.

\subsection{Degree bounds for the Randi\'c index}

\begin{lemma}\label{lem:ch-tree-degree-potential}
For every nonstar tree,
\begin{equation}\label{eq:ch-tree-degree-sum}
 R(T)\ge\sum_{v:\,d_v\ge2}\sqrt{d_v}
                  -\left(\sqrt2-\frac12\right)h.
\end{equation}
Consequently,
\begin{equation}\label{eq:ch-tree-chord}
 R(T)\ge\frac h2+\sqrt\Delta+
                    \frac{L-\Delta}{\sqrt\Delta+\sqrt2}.
\end{equation}
If $\Delta\ge3$, write
\[
 L-2=q(\Delta-2)+b,\qquad 0\le b<\Delta-2,
 \qquad e=\mathbf1_{b>0}.
\]
Then the stronger finite-parameter estimate
\begin{equation}\label{eq:ch-tree-degree-concentrated}
 R(T)\ge\frac h2+q\sqrt\Delta+e\sqrt{b+2}
                                      +(1-q-e)\sqrt2
\end{equation}
also holds.
\end{lemma}

\begin{proof}
Set $\alpha=1/\sqrt2-1/4$, and define
\[
 f(1)=\alpha,\qquad f(d)=\sqrt d-\alpha d\quad(d\ge2).
\]
For an edge with degrees $a,b$, we claim
\[
                  \frac{f(a)}a+\frac{f(b)}b\le\frac1{\sqrt{ab}}.
\]
If one endpoint is a leaf, the other has degree at least two, and this
is an equality. If $a,b\ge2$, put $x=a^{-1/2}$ and $y=b^{-1/2}$.
The assertion is $x+y-xy\le2\alpha$, which follows from
$0<x,y\le1/\sqrt2$: the left side increases in each variable and its
maximum is $\sqrt2-1/2=2\alpha$.
Summing over edges gives $R\ge\sum_v f(d_v)$. Since
\[
 \sum_{d_v\ge2}d_v=2N-L,
 \qquad N-L=h,
\]
this is exactly \eqref{eq:ch-tree-degree-sum}.

The nonleaf degrees satisfy
\[
                   \sum_{d_v\ge2}(d_v-2)=L-2.
\]
The chord inequality on $[2,\Delta]$ gives
\[
 \sqrt d\ge\sqrt2+\frac{d-2}{\sqrt\Delta+\sqrt2}.
\]
Substitute its sum into \eqref{eq:ch-tree-degree-sum}, and use
$\sqrt2+(\Delta-2)/(\sqrt\Delta+\sqrt2)=\sqrt\Delta$ to obtain
\eqref{eq:ch-tree-chord}. The same formula is valid for $\Delta=2$.

For the last estimate, the concavity of $\sqrt x$ shows that, among
$I$ numbers in $[2,\Delta]$ with total excess $L-2$ above two, the
minimum sum of square roots is obtained by placing as much excess as
possible at $\Delta$. Explicitly, if two entries are strictly between
the endpoints, transferring excess from the smaller to the larger until
one reaches an endpoint does not increase the sum. Repetition leaves
$q$ entries equal to $\Delta$, one entry $b+2$ if $b>0$, and all others
equal to two. The resulting minimum is
\[
 q\sqrt\Delta+e\sqrt{b+2}+(I-q-e)\sqrt2.
\]
Substitution proves \eqref{eq:ch-tree-degree-concentrated}.
\end{proof}

\begin{lemma}\label{lem:ch-tree-leaf-bound}
Every nonstar tree satisfies, with $c=1-1/\sqrt2$,
\begin{equation}\label{eq:ch-tree-leaf-lower}
 R(T)\ge\frac h2+\sqrt L+\frac1{\sqrt2}-\frac12-\frac c{\sqrt L}.
\end{equation}
\end{lemma}

\begin{proof}
We induct on the order. Let $v$ be a leaf with neighbor $u$ of degree
$d$. Directly comparing affected edges before and after deleting $v$
gives
\[
 R(T)-R(T-v)=\frac1{\sqrt d}
 -\left(\frac1{\sqrt{d-1}}-\frac1{\sqrt d}\right)
          \sum_{w\in N(u)\setminus\{v\}}\frac1{\sqrt{d_w}}.
\]
If $d=2$, the other neighbor has degree at least two, because $T$ is a
nonstar of order at least four. The difference is therefore at least
$1/2$. The leaf count stays the same and $h$ decreases by one, so this
is exactly the required increment in \eqref{eq:ch-tree-leaf-lower}.

If $d\ge3$, at least one of the other neighbors of $u$ is a nonleaf;
otherwise the connected tree would be a star. Thus the sum in the last
display is at most $d-2+1/\sqrt2=d-1-c$, and the difference is at least
\[
 a(d)=\sqrt d-\sqrt{d-1}
          +c\left(\frac1{\sqrt{d-1}}-\frac1{\sqrt d}\right).
\]
Both summands decrease with $d$. Since $d\le L$, this is at least
$a(L)$. Deletion now reduces $n$ and $L$ by one and leaves $h$ unchanged;
$a(L)$ is exactly the increment of the right side of
\eqref{eq:ch-tree-leaf-lower}.

These comparisons finish the induction whenever $T-v$ is a nonstar.
If $T-v$ is a star, then a nonstar $T$ is obtained by subdividing one arm
of a star, and direct evaluation gives equality in
\eqref{eq:ch-tree-leaf-lower}. This includes the first nonstar, $P_4$,
and establishes the base of the induction.
\end{proof}

\subsection{Exact leaf contraction and distance bounds}

Delete the leaves and give each remaining vertex $v$ weight
$w_v=1+\ell_v$, where $\ell_v$ is its number of original leaf neighbors.
The nonleaves induce a tree $H$ with $h$ edges, and its positive integer
weights total $n$. Let $X,Y$ be independent vertices of $H$ sampled with
probabilities $w_v/n$. The cut formula for tree distances gives the exact
identity
\begin{equation}\label{eq:ch-tree-leaf-contraction}
 \mu(T)=\frac nN\mathbb E d_H(X,Y)+\frac{2L}{n}.
\end{equation}
Indeed, each of the $L$ leaf edges contributes $N$ to the unordered
distance sum. Every other edge contributes precisely the weighted cut
product of its corresponding edge in $H$.

A vertex of degree $\Delta$ has at most $h$ nonleaf neighbors, and hence
its weight is at least $\Delta-h+1$. Also, distributing $L$ leaves
among $I=h+1$ nonleaves yields a weight at least
$1+\lceil L/(h+1)\rceil$. Put
\begin{equation}\label{eq:ch-tree-heavy-a}
 a=\max\left\{\Delta-h+1,\ 1+\left\lceil\frac L{h+1}\right\rceil\right\},
 \qquad \delta=\bigl(2a-(n-h+1)\bigr)_+.
\end{equation}
Lemma~\ref{lem:ch-heavy-tree}, with its nonnegative parity refinement
dropped, and \eqref{eq:ch-tree-leaf-contraction} give
\begin{equation}\label{eq:ch-tree-core-mu}
 \mu(T)\le
 \frac{hn}{2N}-\frac{h(h^2-1)}{6nN}
                   -\frac{h\delta^2}{2nN}+\frac{2L}{n}.
\end{equation}

We will also use the two elementary upper bounds
\begin{align}
 \mu(T)&\le\frac{n+1}{3},\label{eq:ch-tree-path-mu}\\
 \mu(T)&\le2-\frac2n+
          \frac{k(k+1)(3n-2k-7)}{3nN}.
                                     \label{eq:ch-tree-broom-mu}
\end{align}
For \eqref{eq:ch-tree-path-mu}, apply the same weighted-tree lemma
directly to the unweighted tree, with $h=N$ and $a=1$.
For completeness, \eqref{eq:ch-tree-broom-mu} can be proved by starting
with the $\Delta$-star at a maximum-degree vertex and adding the remaining
vertices in an order in which each new vertex is a leaf. At a step whose
resulting order is $t$, sort the new leaf's $t-1$ distances as
$b_1\le\cdots\le b_{t-1}$. A path realizing distance $j$ also realizes
every smaller positive distance, so $b_i\le i$. The current tree still
contains the original $\Delta$-star, and its diameter is at most
$t-\Delta+1$: any path can use at most two of the edges at that center,
and counting the remaining vertices proves the bound. Hence
\[
 \sum_{i=1}^{t-1}b_i\le
 \sum_{i=1}^{t-1}\min\{i,t-\Delta+1\}
 =\binom t2-\binom{\Delta-1}{2}.
\]
The initial star has unordered distance sum $\Delta^2$. Summing these
increments for $t=\Delta+2,\ldots,n$ and substituting
$k=N-\Delta$ proves \eqref{eq:ch-tree-broom-mu}.

\subsection{The infinite range: every order at least 1500}

In \eqref{eq:ch-tree-core-mu}, it is sufficient here to use just
$a\ge\Delta-h+1$. Its correction then has
$\delta\ge(N-2k-h)_+$. Combine this with
\eqref{eq:ch-tree-chord}, use $L=N-h$, and rationalize
$r-s=-k/(s+r)$. The slack satisfies
\begin{equation}\label{eq:ch-tree-master-slack}
\begin{split}
 R-\mu-\gamma_n\ge{}&
 \frac{k-h}{r+\sqrt2}-\frac{k}{s+r}\\
 &+h\left\{\frac{9n+h^2-13}{6nN}
                    +\frac{(N-2k-h)_+^2}{2nN}\right\}.
\end{split}
\end{equation}
Here $h\ge1$ and $k\ge h$. Set
\[
 q=\frac kh\ge1,\qquad z=\frac hN>0,\qquad
 \rho=\sqrt{1-qz}=\frac rs>0,
 \qquad \varepsilon_n=\frac{\sqrt2}s,\qquad C_n=\frac{N\sqrt N}{n}.
\]
In particular $qz<1$. Divide \eqref{eq:ch-tree-master-slack} by $h$
and multiply by $s$. Dropping the positive term
$(9n-13)/(6ns)$ leaves the lower bound
\begin{equation}\label{eq:ch-tree-normalized}
 \frac{q-1}{\rho+\varepsilon_n}-\frac q{1+\rho}
 +C_n\left\{\frac{z^2}{6}
                  +\frac{(1-(2q+1)z)_+^2}{2}\right\}.
\end{equation}
For $n\ge1500$, monotonicity in $n$ and direct squaring give
\[
 \varepsilon_n<\varepsilon_0:=\frac{37}{1000},
 \qquad C_n>C_0:=\frac{3869}{100}.
\]
The exact endpoint checks are
$\varepsilon_0^2>2/1499$ and
$C_0^2<1499^3/1500^2$.

Let
\[
 A(q,z)=\frac{q-1}{\rho+\varepsilon_0}-\frac q{1+\rho},
 \qquad B(q,z)=\frac{z^2}{6}
                          +\frac{(1-(2q+1)z)_+^2}{2}.
\]
Thus it is enough to prove $A+C_0B>0$.
If $q\ge21/10$, the numerator of $A$ after combining fractions is
\[
 q(1-\varepsilon_0)-1-\rho
 \ge\frac{21}{10}(1-\varepsilon_0)-2
 =\frac{223}{10000}>0.
\]
If $z\ge1/2$, use $A\ge-1$, which follows from
\[
 A=(q-1)\left(\frac1{\rho+\varepsilon_0}
                         -\frac1{\rho+1}\right)-\frac1{\rho+1}.
\]
Then $B\ge1/24$, and
$-1+C_0/24=1469/2400>0$.

It remains to cover the fixed compact rectangle
\[
 1\le q\le\frac{21}{10},\qquad 0\le z\le\frac12,
 \qquad qz<1.
\]
Partition the $q$ interval into steps $1/100$ and the $z$ interval into
steps $1/200$. For a rectangle
$[q_-,q_+]\times[z_-,z_+]$ with $q_-z_-<1$, choose rational bounds
\[
 \rho_-\le\sqrt{\max\{0,1-q_+z_+\}},\qquad
 \rho_+\ge\sqrt{1-q_-z_-}
\]
at scale $10^{-6}$. On its feasible part,
\begin{align*}
 A&\ge\frac{q_--1}{\rho_++\varepsilon_0}
                              -\frac{q_+}{1+\rho_-},\\
 B&\ge\frac{z_-^2}{6}
                     +\frac{\max\{0,1-(2q_++1)z_+\}^2}{2}.
\end{align*}
All denominators are positive. The exact certificate
\texttt{ch\_tree\_infinite\_certificate.py} verifies strict positivity
of the sum of these lower bounds on all 10,984 relevant rectangles.
The smallest certified margin is exactly
\[
 \frac{431514416317626190321}{5137900590020000000000}
 >\frac{83}{1000}.
\]
The square-root enclosure is obtained by taking
$j=\lfloor10^6\sqrt x\rfloor$ using integer square roots and using
$j/10^6$ and $(j+1)/10^6$ as endpoints; an exact endpoint may be retained
when its square equals $x$. The certificate checks the squared endpoint
inequalities explicitly. Rectangles with $q_-z_-\ge1$ contain no
feasible point with $qz<1$ and need not be tested.

This is a finite rational verification on a domain independent of $n$.
It proves strict positivity in \eqref{eq:ch-tree-master-slack} for every
nonstar tree of every order $n\ge1500$.

\subsection{The finite range: orders 16 through 1499}

For $4\le n\le1499$, all nonstar trees with $\Delta\ge3$ have parameters
covered by the following loops:
\[
 \begin{gathered}
 1\le h\le n-3,\qquad I=h+1,\qquad L=N-h,\\
 \max\left\{3,\,2+\left\lceil\frac{L-2}{I}\right\rceil\right\}
       \le\Delta\le L.
 \end{gathered}
\]
For each triple, define $q,b,e$ as in
\eqref{eq:ch-tree-degree-concentrated} and retain it when $q+e\le I$.
These conditions follow from the nonleaf degree excess identity and
are therefore necessary for every tree in question.
We compare the maximum of the two lower bounds
\eqref{eq:ch-tree-degree-concentrated} and
\eqref{eq:ch-tree-leaf-lower} with the minimum of the three upper bounds
\eqref{eq:ch-tree-core-mu}, \eqref{eq:ch-tree-path-mu}, and
\eqref{eq:ch-tree-broom-mu}.

Here are the exact arithmetic conventions, which also specify the finite
certificate. Set $S=10^{12}$ and, for positive integers $j\le4000$, put
\[
 \ell_j=\left\lfloor S\sqrt j\right\rfloor,
 \qquad u_j=\ell_j+1,
 \qquad p_j=\left\lfloor\frac{S^2}{u_j}\right\rfloor,
 \qquad t_j=\left\lceil\frac{S^2}{\ell_j}\right\rceil.
\]
These enclose both the square root and its reciprocal:
\[
 \frac{\ell_j}{S}\le\sqrt j<\frac{u_j}{S},\qquad
 \frac{p_j}{S}\le\frac1{\sqrt j}\le\frac{t_j}{S}.
\]
Integer binary search computes $\ell_j$, and the certificate checks
$\ell_j^2\le jS^2<u_j^2$ for every $j$.
The scaled integer lower bounds for $R$ are
\begin{align*}
 B_1&=\frac{hS}{2}+q\ell_\Delta+e\ell_{b+2}
                                      +(1-q-e)u_2,\\
 B_2&=\frac{hS}{2}+\ell_L+p_2-\frac S2-t_L+p_{2L},
 \qquad B=\max\{B_1,B_2\}.
\end{align*}
The $e\ell_{b+2}$ term is omitted when $e=0$.
Here $q\ge1$, since $\Delta\le L$, so the coefficient $1-q-e$ is
nonpositive; its use of the upper endpoint $u_2$ is the required outward
rounding. The $B_2$ formula uses
$c/\sqrt L=1/\sqrt L-1/\sqrt{2L}$.

With $a,\delta$ from \eqref{eq:ch-tree-heavy-a}, set $D_0=6nN$ and
\begin{align*}
 P_{\rm core}&=3hn^2-h(h^2-1)-3h\delta^2+12LN,\\
 P_{\rm broom}&=12N^2+2k(k+1)(3n-2k-7),\\
 P_{\rm path}&=2nN(n+1),\qquad
 P=\min\{P_{\rm core},P_{\rm broom},P_{\rm path}\}.
\end{align*}
The three fractions with common denominator $D_0$ are exactly the three
upper bounds for $\mu$. Therefore the single strict integer comparison
\begin{equation}\label{eq:ch-tree-integer-test}
             (B-u_N+2S)D_0-S(P+12N)>0
\end{equation}
implies $R-\mu-\gamma_n>0$ for every tree with the parameter triple.

The program \texttt{ch\_tree\_finite\_certificate.cpp}, run with
arguments \texttt{16 1499}, checks all 552,444,094 retained triples and
finds no failure of \eqref{eq:ch-tree-integer-test}. The smallest
certified slack is approximately $0.005921876485$, attained among these
bounds at $(n,h,L,\Delta)=(16,9,6,5)$. This displayed decimal is only a
report of an exact rational comparison; no floating-point number is
used to accept a case or choose the minimum.

All proof decisions in that program use signed 128-bit integers, with
64-bit storage for individual scaled bounds and polynomial numerators.
There is no overflow in the stated range. For example, $D_0<1.4\cdot10^7$,
the absolute values of the three $P$ numerators are below $3\cdot10^{10}$,
and $|B|<4000S$. The latter follows from
$q\Delta+e(b+2)\le2N-L<3000$, together with the displayed formulas.
Thus the absolute value of the integer margin in
\eqref{eq:ch-tree-integer-test} is below $9\cdot10^{22}$, and its
cross-products for comparing margins with different denominators are
below $1.3\cdot10^{30}$. The binary-search root squares are below
$10^{28}$. In particular every 128-bit intermediate has absolute value
less than $10^{31}<2^{127}$, while every 64-bit intermediate is smaller
than $10^{17}<2^{63}$.

\subsection{The seven small profiles and the paths}

The independent rational script
\texttt{ch\_tree\_small\_certificate.py} first checks the same scalar
bounds for $4\le n\le15$. There are 239 retained parameter triples.
Exactly 232 satisfy the strict scalar comparison, and the seven remaining
triples are listed below. The failure of a scalar comparison only means
that those particular bounds leave that profile undecided.

\[
\begin{array}{c|c|c|c|r}
 n&L&\Delta&n-L-2&\text{ordered subdivision choices}\\\hline
 11&5&4&4&126\\
 12&6&5&4&210\\
 12&5&4&5&252\\
 13&6&5&5&462\\
 13&5&4&6&462\\
 14&6&5&6&924\\
 15&6&5&7&1716
\end{array}
\]

Every tree in one of these seven profiles has exactly two vertices of
degree at least three, of degrees $(\Delta,3)$. Indeed, for $(L,\Delta)$
equal to $(5,4)$ or $(6,5)$, the maximum-degree vertex consumes all but
one of the total excess $L-2$ above degree two. Suppressing degree-two
vertices therefore leaves two adjacent hubs with $\Delta-1$ pendant
edges at the first and two at the second. There are $L+1$ edges in this
reduced tree. Every original tree is obtained by distributing its
$n-L-2$ degree-two vertices among these $L+1$ edges.

The script enumerates every ordered weak composition of $n-L-2$ into
$L+1$ parts. Their number is $\binom{n-2}{L}$ for each row, giving 4,152
choices in total. Permutations of equivalent pendant arms may repeat a
tree, which does not affect coverage. For each generated tree it computes
the unordered distance sum exactly as
\[
                      W(T)=\sum_{e\in E(T)}s_e(n-s_e)
\]
using integer component orders $s_e$. It bounds each Randi\'c edge term
below by $1/\operatorname{upper}(\sqrt{d_ud_v})$ at scale $10^{-12}$,
and bounds $\sqrt N$ above at the same scale. Every resulting rational
lower bound for $R-2W/(nN)-\gamma_n$ is strictly positive. The smallest
displayed bound is approximately $0.11502894422983344$.
This covers all seven profiles without requiring an external list of
unlabeled trees.

Finally, trees with $\Delta=2$ are paths and were omitted from the scalar
loops. For $n\ge3$ their exact slack is
\[
 f(n)=R(P_n)-\mu(P_n)-\gamma_n
      =\frac{n+1}{6}+\sqrt2-\sqrt{n-1}-\frac2n.
\]
For real $x\ge10$,
\[
 f'(x)=\frac16-\frac1{2\sqrt{x-1}}+\frac2{x^2}>0,
 \qquad f(10)=\sqrt2-\frac{41}{30}>0.
\]
For $4\le n\le9$, a square-root enclosure at scale $10^{-12}$ gives
the following exact positive lower bounds:
\[
\begin{array}{c|c}
 n&\text{rational lower bound for }f(n)\\\hline
 4&11622066103/750000000000\\
 5&3553390593/250000000000\\
 6&34436754619/3000000000000\\
 7&259200211369/21000000000000\\
 8&4615562827/250000000000\\
 9&136038969317/4500000000000
\end{array}
\]
Each is greater than $11/1000$. These entries are obtained from the
explicit expression
\[
 \frac{n+1}{6}
 +\frac{\lfloor10^{12}\sqrt2\rfloor}{10^{12}}
 -\frac{\lfloor10^{12}\sqrt{n-1}\rfloor+1}{10^{12}}
 -\frac2n,
\]
so they can be checked using integer squares alone. The paths of orders
two and three are stars and have already been included in the equality
case.

The analytic infinite-range argument, the finite parameter certificate,
the seven-profile enumeration, and the path computation together exhaust
every nonstar tree of every order. Every such tree has strict inequality,
and the stars attain equality. This proves
Theorem~\ref{thm:ch-all-trees}.

\paragraph{Reproducible certificate sources.}
Appendix A.1--A.3 supplies the three complete source files.
The infinite-order and small-order certificates use Python integer
arithmetic and \texttt{Fraction}; the finite-order C++ certificate uses
integer square-root enclosures and signed \texttt{\_\_int128} comparisons.
Decimal output is descriptive only. The finite ranges and parameter
coverage are proved above, rather than assumed from numerical sampling.


\section{The additive inequality for every graph without leaves}
\label{sec:ch-min-degree-two}

This section removes the order restriction entirely when the minimum
degree is at least two. Its only corpus input is
\[
 M\ge\frac{(n-1)\mu-(n-2)}n.\tag{U}
\]
In particular the finite verification at the end concerns integer
parameters, and is exhaustive over all parameters not covered by the
analytic argument. It is not an enumeration of selected graphs.

Write \(\mathcal W(G)=\sum_{\{u,v\}}d_G(u,v)=\binom n2\mu(G)\).
For \(d\le n-1\), define the broom mean-distance bound
\[
 B(n,d)=2-\frac2n+
 \frac{k(k+1)(3n-2k-7)}{3n(n-1)},\qquad k=n-1-d.
\]
The corresponding broom consists of a vertex \(v\) with \(d\) incident
arms, \(d-1\) of length one and one of length \(k+1\).

\begin{lemma}[A strict broom improvement without leaves]
\label{lem:ch-leafless-broom}
If \(G\) is a connected simple graph of order \(n\ge3\), minimum degree
at least two and maximum degree \(d\), then
\[
 \mu(G)\le B(n,d)-\frac{k}{\binom n2},
 \qquad k=n-1-d.
\]
\end{lemma}

\begin{proof}
Choose a maximum-degree vertex \(v\) and a spanning tree \(T\) containing
all its \(d\) incident edges. We first treat \(d\ge3\).

Build \(T\) from that initial \(d\)-star by adding its other vertices as
leaves. At a step whose new order is \(t\), the new leaf's ordered
distances \(b_1\le\cdots\le b_{t-1}\) satisfy
\[
 b_i\le\min\{i,t-d+1\}.
\]
The first bound follows from a path to a farthest vertex; the second is
the diameter bound for a tree containing the original \(d\)-star.
The new transmission is consequently at most
\[
 U(t)=\binom t2-\binom{d-1}{2}.
\]
Equality here forces a single vertex at each distance
\(1,\ldots,t-d\) and \(d-1\) vertices at distance \(t-d+1\).
The tree is therefore a broom; since \(d\ge3\), its unique vertex of
degree \(d\) is the original \(v\).

Suppose first that \(T\) is not a broom centered at \(v\).
At the first step that leaves the broom family, the previous tree has
order \(d+L\) and long-arm length \(L\ge2\), and the new order is
\(d+L+1\). Attaching another vertex at \(v\) is impossible because its
full degree \(d\) is already present. Attaching at the long-arm terminal
would preserve the broom family. The remaining possibilities have the
following transmission losses relative to attachment at that terminal:
\[
 \begin{array}{c|c}
 \text{attachment point in the previous broom}&\text{loss}\\ \hline
 \text{one of the other root leaves}&(d-2)(L-1)\\
 \text{long-arm vertex at distance }j,\ 1\le j\le L-1&
       (L-j)(d+j-1).
 \end{array}
\]
These formulas follow by subtracting the two vertex transmissions in
the previous broom. Each loss is at least \(L-1\). In the second case
the quadratic in \(j\) is concave, so its minimum is at an endpoint and
equals at least
\(\min\{d(L-1),d+L-2\}\ge L-1\).

A nonbroom shape centered at this fixed \(v\) cannot become a broom by
subsequent leaf additions: either two arms are already nontrivial, or
some vertex away from \(v\) already has degree at least three. Equality
in the transmission bound at any later step is therefore impossible.
All distances are integers, so each of the remaining \(k-L\) additions
loses at least one. Summing the losses gives
\[
 \mathcal W(T)\le\binom n2 B(n,d)-(L-1)-(k-L)
 =\binom n2B(n,d)-(k-1).
\]
Since the minimum degree of \(G\) is at least two, \(G\) has at least
\(n\) edges. There is an edge of \(G\) outside \(T\); its endpoints have
distance at least two in \(T\), and distance one in \(G\). All other
distances can only decrease. Hence
\[
 \mathcal W(G)\le\mathcal W(T)-1
 \le\binom n2B(n,d)-k.
\]

Suppose instead that \(T\) is the broom. Let \(w\) be its long-arm
terminal. This arm has length \(k+1\). Because \(d_G(w)\ge2\), there is
an edge \(wz\) of \(G\) outside \(T\). If \(z\) is another root leaf,
the distance of this single pair drops from \(k+2\) to one, saving
\(k+1\). Otherwise \(z\) lies on the long arm at distance \(j\ge2\)
from \(w\). For every vertex on the root side of \(z\), the chord saves
at least \(j-1\) in its distance from \(w\). There are \(n-j\) such
vertices. The saving is at least
\[
 (n-j)(j-1)\ge
 \min\{n-2,dk\}\ge k,\qquad 2\le j\le k+1.
\]
The first inequality follows from concavity in \(j\). This proves the
same Wiener-index improvement. When \(k=0\), the required improvement
is zero and the original broom bound suffices.

Finally, if \(d=2\), connectedness and minimum degree two force \(G\)
to be a cycle. Delete an appropriate cycle edge to obtain a spanning
path whose chosen degree-two vertex is adjacent to one endpoint.
This is the \(d=2\) broom, and the preceding long-terminal chord
argument applies verbatim. Dividing by \(\binom n2\) finishes the proof.
\end{proof}

\begin{theorem}[All orders, minimum degree at least two]
\label{thm:ch-min-degree-two}
Let \(G\) be a connected simple graph of order \(n\ge3\) with minimum
degree at least two. If \(G\) satisfies \textup{(U)}, then
\[
 R(G)-\mu(G)>\sqrt{n-1}-2+\frac2n.
\]
\end{theorem}

\begin{proof}
Set
\[
 s=\sqrt{n-1},\quad d=\Delta(G),\quad r=\sqrt{d/2},
 \quad\gamma=s-2+\frac2n,\quad X=R-M.
\]
Assume, for a contradiction, that \(g=R-\mu\le\gamma\).
The uniform input gives
\[
 0\le X\le g+1+\frac{\mu-2}{n}
       \le s-1+\frac\mu n.
\]
For an edge whose endpoint degrees satisfy \(2\le a\le b\le d\),
\[
 (a^{-1/2}-b^{-1/2})^2
 =(\sqrt{b/a}-1)((ab)^{-1/2}-b^{-1})
 \le(r-1)((ab)^{-1/2}-b^{-1}).
\]
Summing proves the sharper degree-ratio bridge
\[
 n-2R\le(r-1)(R-M).
\]
It follows that
\begin{equation}\label{eq:ch-delta2-necessary}
 n\le 2(\mu+\gamma)+(r-1)\left(s-1+\frac\mu n\right).
\end{equation}
The right side increases with \(\mu\).

\paragraph{All \(n\ge256\).}
Here the original broom bound \(B=B(n,d)\) already suffices. Put
\[
 c=\frac dn,\quad \varepsilon=\frac1n,\quad \beta=\frac Bn,
 \quad b(c)=\frac{1-3c^2+2c^3}{3}.
\]
An exact expansion gives
\[
 \beta=b(c)+\frac{\varepsilon}{1-\varepsilon}
 \left[
  \frac{1+9c-6c^2+2c^3}{3}
  -\frac{7+5c}{3}\varepsilon+2\varepsilon^2
 \right]
 \le b(c)+\frac2{n-1}.
\]
Indeed the first bracketed polynomial increases from \(1/3\) to \(2\)
on \([0,1]\), and the remaining correction is nonpositive.
Also \(0<\beta<1\), since \(B\le(n+1)/3\).

After substituting \(\mu\le B\) in
\eqref{eq:ch-delta2-necessary}, its right side divided by \(n\) becomes
\[
 2\beta+\sqrt{\frac c2}\sqrt{1-\varepsilon}
 +\sqrt\varepsilon\sqrt{1-\varepsilon}
 +\sqrt{\frac{c\varepsilon}{2}}(\beta-1)
 -(3+\beta)\varepsilon+4\varepsilon^2.
\]
The last two displayed correction terms are nonpositive for \(n\ge3\).
Consequently the necessary inequality implies
\begin{equation}\label{eq:ch-delta2-asymptotic}
 1\le2b(c)+\sqrt{\frac c2}+\frac4{n-1}+\frac1{\sqrt n}.
\end{equation}

The following elementary uniform estimate has adequate slack:
\[
 2b(c)+\sqrt{c/2}<\frac{23}{25},
 \qquad 0\le c\le1.
\]
To verify it, concavity at \(c=1/4\) gives
\[
 \sqrt{c/2}\le\frac{c+1/4}{\sqrt2}
 <\frac{71}{100}(c+1/4).
\]
Put \(a=71/100\) and
\(p(c)=2/3-2c^2+4c^3/3+a(c+1/4)\).
Its only interior local maximum occurs at
\[
 c_0=\frac{1-\sqrt{1-a}}2<\frac{231}{1000}.
\]
At a critical point \(a=4c_0(1-c_0)\), so
\[
 p(c_0)=\frac23+c_0+c_0^2-\frac83c_0^3
 <\frac23+\frac{231}{1000}
   +\left(\frac{231}{1000}\right)^2
   -\frac83\left(\frac{231}{1000}\right)^3
 <\frac{23}{25}.
\]
The intervening cubic expression increases on \([0,1/4]\); both
endpoint values \(p(0)\) and \(p(1)\) are also below \(23/25\).
Thus \eqref{eq:ch-delta2-asymptotic} contradicts the rational inequality
\[
 \frac{23}{25}+\frac4{255}+\frac1{16}
 =1-\frac{37}{20400}<1.
\]

\paragraph{The remaining finite parameter range.}
For \(3\le n\le255\), use Lemma~\ref{lem:ch-leafless-broom} and put
\[
 B_2=B(n,d)-\frac{2(n-1-d)}{n(n-1)}.
\]
For every pair of integers
\[
 3\le n\le255,\qquad 2\le d\le n-1,
\]
the exact inequality
\begin{equation}\label{eq:ch-delta2-finite}
 2(B_2+\gamma)+(r-1)(s-1+B_2/n)<n
\end{equation}
holds. There are \(32131\) pairs. The following self-contained rational
certificate verifies every pair. Each radical is replaced by a rational
upper bound, which can only enlarge the left side. No floating-point
arithmetic is used.

\begin{lstlisting}[basicstyle=\ttfamily\footnotesize,breaklines=true]
from fractions import Fraction as F
from math import isqrt
Q = 10**6
def sqrt_upper(x):
    x = F(x)
    y = x*Q*Q
    a = isqrt(y.numerator // y.denominator)
    if F(a*a,Q*Q) < x:
        a += 1
    return F(a,Q)
count = 0
minimum = None
for n in range(3,256):
    s = sqrt_upper(n-1)
    gamma = s-2+F(2,n)
    for d in range(2,n):
        k = n-1-d
        B2 = (2-F(2,n)
              +F(k*(k+1)*(3*n-2*k-7),3*n*(n-1))
              -F(2*k,n*(n-1)))
        r = sqrt_upper(F(d,2))
        margin = n-(2*(B2+gamma)+(r-1)*(s-1+B2/n))
        assert margin > 0
        minimum = margin if minimum is None else min(minimum,margin)
        count += 1
assert count == 32131
assert minimum == F(22956105475381,363687500000000)
\end{lstlisting}
The least certified margin occurs at \((n,d)=(23,6)\) and exceeds
\(0.06312\). Inequality~\eqref{eq:ch-delta2-finite} contradicts
\eqref{eq:ch-delta2-necessary}, exhausting the finite range. This proves
the strict inequality for every order.
\end{proof}

\paragraph{Consequences for a remaining counterexample.}
Together with the general-order cutoff proved below, this result forces every
remaining counterexample to the additive inequality to have at least
one leaf and order below \(3\,000\,000\). The proof does not establish
that deleting such a leaf preserves the required defect inequality;
that is a separate structural question.


\section{A maximum-degree-aware cutoff of three million}
\label{sec:ch-three-million}

The argument in this section assumes only the uniform corpus input
\[
 M\ge\frac{(n-1)\mu-(n-2)}n.\tag{U}
\]
The graph is finite, simple and connected; $R$, $M$ and $\mu$ retain their
preceding definitions. All remaining estimates are elementary. The finite
certificate below checks rational parameter inequalities, not a sample of
graphs.



\begin{theorem}\label{thm:ch-three-million}
For every $n\ge3\,000\,000$, a connected simple graph satisfying
\textup{(U)} satisfies
\[
 R-\mu\ge\sqrt{n-1}-2+\frac2n,
\]
with equality precisely for $K_{1,n-1}$.
\end{theorem}

\begin{proof}
Write
\[
 s=\sqrt{n-1},\quad r=\sqrt\Delta,\quad
 k=n-1-\Delta,\quad \gamma=s-2+\frac2n,
 \quad g=R-\mu,\quad X=R-M.
\]
Assume $g\le\gamma$ and show that the graph is a star. From
\textup{(U)} and $M\le R$,
\[
 0\le X\le g+1+\frac{\mu-2}{n}
       \le s-1+\frac\mu n<s<\sqrt n.
\]
The elementary edge identity gives
\[
 n-2R\le(r-1)X.
\]
Consequently, on putting $\eta=\Delta/n$,
\begin{equation}\label{eq:heavy-mu-lower}
 \frac\mu n>
 \frac{1-\sqrt\eta}{2}-\frac1{2\sqrt n}.
\end{equation}
Indeed $n-2\mu-2\gamma\le(r-1)X<(r-1)s$, and the resulting lower bound
is at least the right side of \eqref{eq:heavy-mu-lower}.

\paragraph{The two outer degree ranges.}
If $\eta\le1/10$, the path bound $\mu\le(n+1)/3$ contradicts
\eqref{eq:heavy-mu-lower}; the exact rational comparison is
\[
 \frac{1-31623/100000}{2}-\frac1{3464}
 >\frac{1+1/3000000}{3}.
\]
Here $(31623/100000)^2>1/10$ and $1732^2<3\,000\,000$.

If $\eta\ge7/10$ and $k\ge1$, combine the edge identity,
\textup{(U)}, and the elementary broom distance bound
\[
 \mu\le2-\frac2n+
 \frac{k(k+1)(3n-2k-7)}{3n(n-1)}.
\]
Combining these inequalities gives, with $A_r=2+(r-1)/n$,
\begin{equation}\label{eq:heavy-final-forced}
 \frac{s-1}{s+r}\le
 \frac{2(r-1)(n-1)}{n^2k}+
 A_r\frac{(k+1)(3n-2k-7)}{3n(n-1)}.
\end{equation}
For clarity, before substituting the broom bound that algebra reads
\[
 (r+1)g\ge n-2\mu-(r-1)\left(1+\frac{\mu-2}{n}\right),
 \qquad
 (s-1)(s-r)\le\frac{2(r+1)}n+A_r(\mu-2).
\]
Use $s-r=k/(s+r)$ and divide by $k$ to get
\eqref{eq:heavy-final-forced}. Set $x=k/n<3/10$. Its final fraction is at
most $x-2x^2/3+2/n\le6/25+2/n$: writing $\varepsilon=1/n$ bounds it by
$(x-2x^2/3+\varepsilon)/(1-\varepsilon)$, and
$x-2x^2/3+2\varepsilon<1$ justifies the stated estimate.
The two sides of \eqref{eq:heavy-final-forced} are incompatible because
\[
 \frac12-\frac1{3464}
 >\frac2{1732}+
  \left(2+\frac1{1732}\right)
  \left(\frac6{25}+\frac1{1500000}\right).
\]
The case $k=0$ is handled at the end. It remains to exclude
$1/10<\eta<7/10$.

\paragraph{Packing that retains the maximum-degree vertex.}
Set $K=\lceil7n^{1/3}\rceil$ and $t=\sqrt{K-1}$.
In the remaining degree range $K\le\Delta$. Choose a maximal collection
of $q$ vertices of degree at least $K$, at pairwise distance at least
three, starting with a maximum-degree vertex. Their closed radius-one
balls are disjoint, so
\[
 \Delta+1+(q-1)(K+1)\le n,
 \qquad
 q\le1+\frac{(1-\eta)n^{2/3}}7.\tag{P}
\]
Let $A$ be the union of their closed radius-three balls, and let
$b=n-|A|$. Every edge incident to a vertex of degree at least $K$ has
both endpoints in $A$. Splitting the degree-difference identity across
those edges and the other edges therefore gives
\[
 R\ge\frac b2-\frac{t-1}{2}X.\tag{S}
\]
This follows because the first group contributes at most $|A|$ to
$n-2R$, whereas the other edges have degree ratio at most $K-1$ and
contribute at most $(t-1)X$.

Contract the components of $G[A]$, give each its original order as weight,
and take a spanning tree of the quotient. Let it have $h$ edges. The
number of contracted components is at most $q$, and hence
$b\ge h-q+1$. One contracted vertex has weight at least $\Delta+1$;
in particular $h\le k$. A component containing $q_i$ of the selected
balls has internal diameter at most $7q_i-1$. A path in the quotient
tree lifts with at most $7q$ extra edges. Thus Lemma~\ref{lem:ch-heavy-tree}
applies with $a=\Delta+1$.

Put $z=h/n$, and define, for $y\ge0$,
\[
 F(y,c)=y^3+3y(2c+y-1)_+^2,
 \qquad \Phi(y,c)=\frac y2-\frac{F(y,c)}6.
\]
The heavy-vertex correction is at least
$h(2\eta+z-1)_+^2/2$. Therefore
\begin{equation}\label{eq:heavy-mu-upper}
 \frac\mu n\le\frac1{1-1/n}
 \left(\Phi(z,\eta)+\frac1{6n^2}+\frac7n
                  +(1-\eta)n^{-1/3}\right).
\end{equation}
Here the term $z/(6n^2)$ was enlarged to $1/(6n^2)$ and \textup{(P)}
was used. Notice that $0\le z\le1-\eta$.

Subtract the lifted distance bound from \textup{(S)} and use
$b\ge h-q+1$. The cancellation $h/[2(n-1)]-1/2\le0$ gives
\[
 \frac{h^3}{6n(n-1)}+
 \frac{h n(2\eta+z-1)_+^2}{2(n-1)}
 \le g+\frac{t-1}{2}X+
 \left(\frac12+\frac{7n}{n-1}\right)q+\frac1{6n}.
\]
Since $g<\sqrt n$, $X<\sqrt n$ and $t<\sqrt7\,n^{1/6}$,
multiplication by $6(n-1)/n^2$ yields
\begin{equation}\label{eq:heavy-cubic}
 F(z,\eta)\le
 \left(3\sqrt7+\frac{(45+42/(n-1))(1-\eta)}7\right)n^{-1/3}
 +3n^{-1/2}+\frac{45+42/(n-1)}n+\frac1{n^2}.
\end{equation}

\paragraph{A finite rational interval certificate.}
Put
\[
 e=\frac1{3000000},\quad u=\frac1{144},\quad v=\frac1{1732},
 \quad a_0=\frac{9000003}{200000},\quad b_0=\frac{2646}{1000},
\]
and
\[
 E(c)=\left(3b_0+\frac{a_0(1-c)}7\right)u+3v+a_0e+e^2.
\]
For all $n\ge3\,000\,000$ these are valid outward bounds:
$n^{-1}\le e$, $n^{-1/3}<u$, $n^{-1/2}<v$,
$\sqrt7<b_0$ and $45+42/(n-1)<a_0$.

Partition $[1/10,7/10]$ into the sixty intervals
$[c_-,c_+]=[j/100,(j+1)/100]$, $10\le j\le69$.
For each interval the following table gives a rational decimal $H$ such
that
\[
 0<H\le1-c_-,\qquad F(H,c_-)>E(c_-).
\]
The function $F(y,c)$ increases with both arguments. Thus
\eqref{eq:heavy-cubic} forces $z<H$. Also $\Phi(y,c)$ decreases in $c$
and increases in $y$ for $0\le y\le1-c$. Consequently
\eqref{eq:heavy-mu-upper} bounds $\mu/n$ by
\[
 U_-=\frac{\Phi(H,c_-)+e^2/6+7e+(1-c_-)u}{1-e}.
\]
On the other hand, \eqref{eq:heavy-mu-lower} requires $\mu/n>L_+$, where
\[
 L_+=\frac{1-S_+}{2}-\frac v2,
 \qquad
 S_+=10^{-5}\left\lceil10^5\sqrt{c_+}\right\rceil.
\]
Every row has $U_-<L_+$.
All comparisons use integers or fractions; the displayed upper bounds
are rounded upward and displayed lower bounds downward to six decimal
places. The smallest exact gap exceeds $0.0027308$.

% Embed the contents of ch_heavy_cutoff_table.tex at this location in a final
% single-file manuscript; the separate file is an internal assembly component.
% Generated by ch_heavy_cutoff_certificate.py; displayed decimals are rounded.
\begin{longtable}{rrrrr}
\toprule
$c_-$ & $c_+$ & $H$ & upper bound for $\mu/n$ & lower bound for $\mu/n$\\
\midrule\endhead
0.10 & 0.11 & 0.45956 & 0.219857 & 0.333876\\
0.11 & 0.12 & 0.45885 & 0.219507 & 0.326501\\
0.12 & 0.13 & 0.45814 & 0.219157 & 0.319431\\
0.13 & 0.14 & 0.45743 & 0.218807 & 0.312626\\
0.14 & 0.15 & 0.45672 & 0.218457 & 0.306061\\
0.15 & 0.16 & 0.45600 & 0.218103 & 0.299711\\
0.16 & 0.17 & 0.45529 & 0.217752 & 0.293551\\
0.17 & 0.18 & 0.45457 & 0.217397 & 0.287576\\
0.18 & 0.19 & 0.45385 & 0.217042 & 0.281766\\
0.19 & 0.20 & 0.45312 & 0.216682 & 0.276101\\
0.20 & 0.21 & 0.45240 & 0.216327 & 0.270581\\
0.21 & 0.22 & 0.45167 & 0.215967 & 0.265186\\
0.22 & 0.23 & 0.45094 & 0.215607 & 0.259916\\
0.23 & 0.24 & 0.45021 & 0.215246 & 0.254761\\
0.24 & 0.25 & 0.44947 & 0.214882 & 0.249711\\
0.25 & 0.26 & 0.44873 & 0.214517 & 0.244756\\
0.26 & 0.27 & 0.44799 & 0.214152 & 0.239901\\
0.27 & 0.28 & 0.44725 & 0.213787 & 0.235131\\
0.28 & 0.29 & 0.44641 & 0.213372 & 0.230451\\
0.29 & 0.30 & 0.44442 & 0.212381 & 0.225846\\
0.30 & 0.31 & 0.44119 & 0.210772 & 0.221321\\
0.31 & 0.32 & 0.43695 & 0.208657 & 0.216866\\
0.32 & 0.33 & 0.43186 & 0.206116 & 0.212481\\
0.33 & 0.34 & 0.42603 & 0.203207 & 0.208161\\
0.34 & 0.35 & 0.41957 & 0.199981 & 0.203906\\
0.35 & 0.36 & 0.41255 & 0.196476 & 0.199711\\
0.36 & 0.37 & 0.40504 & 0.192726 & 0.195571\\
0.37 & 0.38 & 0.39709 & 0.188756 & 0.191486\\
0.38 & 0.39 & 0.38875 & 0.184591 & 0.187461\\
0.39 & 0.40 & 0.38006 & 0.180251 & 0.183481\\
0.40 & 0.41 & 0.37106 & 0.175756 & 0.179551\\
0.41 & 0.42 & 0.36178 & 0.171121 & 0.175671\\
0.42 & 0.43 & 0.35226 & 0.166365 & 0.171836\\
0.43 & 0.44 & 0.34251 & 0.161496 & 0.168046\\
0.44 & 0.45 & 0.33258 & 0.156536 & 0.164296\\
0.45 & 0.46 & 0.32249 & 0.151496 & 0.160591\\
0.46 & 0.47 & 0.31226 & 0.146386 & 0.156926\\
0.47 & 0.48 & 0.30193 & 0.141225 & 0.153296\\
0.48 & 0.49 & 0.29151 & 0.136020 & 0.149711\\
0.49 & 0.50 & 0.28103 & 0.130786 & 0.146156\\
0.50 & 0.51 & 0.27053 & 0.125541 & 0.142636\\
0.51 & 0.52 & 0.26003 & 0.120295 & 0.139151\\
0.52 & 0.53 & 0.24955 & 0.115060 & 0.135701\\
0.53 & 0.54 & 0.23912 & 0.109851 & 0.132286\\
0.54 & 0.55 & 0.22878 & 0.104685 & 0.128901\\
0.55 & 0.56 & 0.21854 & 0.099571 & 0.125541\\
0.56 & 0.57 & 0.20845 & 0.094530 & 0.122216\\
0.57 & 0.58 & 0.19852 & 0.089570 & 0.118921\\
0.58 & 0.59 & 0.18879 & 0.084709 & 0.115651\\
0.59 & 0.60 & 0.17927 & 0.079955 & 0.112411\\
0.60 & 0.61 & 0.17000 & 0.075325 & 0.109196\\
0.61 & 0.62 & 0.16100 & 0.070830 & 0.106006\\
0.62 & 0.63 & 0.15229 & 0.066480 & 0.102846\\
0.63 & 0.64 & 0.14389 & 0.062285 & 0.099711\\
0.64 & 0.65 & 0.13581 & 0.058250 & 0.096596\\
0.65 & 0.66 & 0.12807 & 0.054384 & 0.093506\\
0.66 & 0.67 & 0.12067 & 0.050690 & 0.090441\\
0.67 & 0.68 & 0.11362 & 0.047170 & 0.087396\\
0.68 & 0.69 & 0.10693 & 0.043830 & 0.084376\\
0.69 & 0.70 & 0.10059 & 0.040665 & 0.081376\\
\bottomrule
\end{longtable}


These sixty contradictions exhaust the middle degree range.

\paragraph{The remaining dominating-vertex case.}
We have proved that $g\le\gamma$ forces $k=0$. The edges incident to a
dominating vertex contribute one to $M$, so $M\ge1$. The edge bridge
with $r=s$ gives
\[
 R\ge\frac{n+(s-1)M}{s+1}\ge\sqrt{n-1}.
\]
All pair distances are at most two, whence
$\mu=2-|E(G)|/\binom n2\le2-2/n$. Any edge in addition to the spanning
star makes this last inequality strict. Equality is therefore attained
precisely by $K_{1,n-1}$, as claimed.
\end{proof}

\paragraph{Scope.}
This lowers the sufficient cutoff from $10^{10}$ to $3\,000\,000$ with
the same input \textup{(U)}. It leaves $10\le n<3\,000\,000$ beyond
the previously recorded small-order verification. The argument neither
claims a new validation of \textup{(U)} nor settles the remaining
finite interval.


\section{A least counterexample cannot have a subdivided pendant edge}
\label{sec:ch-leaf-support-two}

The following argument uses neither the corpus weighted theorem nor a
conjectural deletion principle. Put
\[
 \gamma_n=\sqrt{n-1}-2+\frac2n,\qquad g(G)=R(G)-\mu(G).
\]

\begin{lemma}[A universal leaf-deletion distance bound]
\label{lem:ch-leaf-distance-increment}
Let \(v\) be a leaf of a connected graph \(G\) of order \(n\ge3\). Then
\[
 \mu(G)-\mu(G-v)\le
 \begin{cases}
  \displaystyle\frac5{12}+\frac1{4(n-1)},&n\text{ even},\\[4pt]
  \displaystyle\frac5{12}+\frac1{4(n-2)},&n\text{ odd}.
 \end{cases}
\]
\end{lemma}

\begin{proof}
Let \(\tau_v=\sum_xd_G(v,x)\), and let
\(t_j=|\{x:d_G(v,x)\ge j\}|\), \(1\le j\le\operatorname{ecc}(v)\).
Every distance level up to the eccentricity is nonempty, so the \(t_j\)
are distinct, strictly decreasing integers in \(\{1,\ldots,n-1\}\).
The elementary identities and inequality
\[
 \tau_v=\sum_jt_j,\qquad
 \mathcal W(G)\ge
 \sum_{\{x,y\}}|d(v,x)-d(v,y)|
 =\sum_j t_j(n-t_j)
\]
give
\[
 \tau_v-(n-1)\mu(G)
 =\tau_v-\frac{2\mathcal W(G)}n
 \le\sum_jt_j\left(\frac{2t_j}n-1\right)
 \le\sum_{t=\lceil n/2\rceil}^{n-1}
             t\left(\frac{2t}n-1\right).
\]
In the last step we discarded negative terms and included every possible
positive term; distinctness of the \(t_j\) is essential.

Deletion of a leaf preserves every distance among the remaining
vertices. Thus
\[
 \mu(G)-\mu(G-v)
 =\frac{2\{\tau_v-(n-1)\mu(G)\}}{(n-1)(n-2)}.
\]
For \(n=2m\), the displayed sum is
\((5m^2-6m+1)/6\). For \(n=2m+1\), it is
\(m(5m-1)/6\). Substitution gives the claimed bounds.
\end{proof}

\begin{theorem}[Degree-two leaf supports are reducible above order 43]
\label{thm:ch-leaf-support-two}
Suppose \(G\) is connected, has order \(n\ge44\), and has a leaf \(v\)
whose neighbor has degree two. If \(G-v\) satisfies the Caporossi--Hansen
additive inequality, then \(G\) satisfies it strictly.
\end{theorem}

\begin{proof}
Write \(u\) for the neighbor of \(v\) and \(w\) for the other neighbor of
\(u\). Connectedness and \(n\ge44\) imply \(d_G(w)\ge2\). Only the two
edges incident to \(u\) change their Randi\'c contributions, giving
\[
 R(G)-R(G-v)
 =\frac1{\sqrt2}
  -\frac{1-1/\sqrt2}{\sqrt{d_G(w)}}\ge\frac12.
\]
Lemma~\ref{lem:ch-leaf-distance-increment} therefore implies
\[
 g(G)-g(G-v)
 \ge\frac1{12}-\frac1{4(n-2)}.
\]
Set \(y=\sqrt{n-2}\). For \(n\ge44\),
\[
 y\ge\sqrt{42}>3+2\sqrt3,
 \qquad y^2-6y-3>0.
\]
Consequently
\[
 \frac1{12}-\frac1{4(n-2)}
 >\frac1{2\sqrt{n-2}}
 >\gamma_n-\gamma_{n-1},
\]
where
\[
 \gamma_n-\gamma_{n-1}
 =\frac1{\sqrt{n-1}+\sqrt{n-2}}
  -\frac2{n(n-1)}.
\]
If \(g(G-v)\ge\gamma_{n-1}\), addition yields \(g(G)>\gamma_n\).
\end{proof}

Thus a least-order counterexample of order at least \(44\), if one
exists, has no leaf adjacent to a degree-two vertex. In combination with
the all-order minimum-degree-two result, a remaining least counterexample
above this order must have leaves, and every leaf must be adjacent to a
vertex of degree at least three. This conclusion is a structural
restriction, not a proof that such a graph cannot exist.


\section{What remains for the unrestricted conjecture}
\label{sec:new-remaining}

The three main results do not partition all connected graphs into solved
classes. Their uncovered intersection is precise: a graph must have both
at least one cycle and at least one leaf, and have order below
$3\,000\,000$. Relative to the verification through order nine recorded
in [227], a counterexample would therefore satisfy
\[
 10\le n<3\,000\,000,\qquad
 |E(G)|\ge n,\qquad \delta(G)=1.
\]
If one chooses a counterexample of least order and that order is at
least $44$, every leaf has a neighbor of degree at least three.
This last condition concerns a \emph{least-order} counterexample, not
every graph that is currently outside the proved classes.

\subsection{Why the present scalar bounds do not finish the mixed case}

Here is an exact feasible point for the scalar relaxation used in our
mixed-graph investigations:
\[
\begin{gathered}
 n=1000,\quad L=447,\quad I=553,\quad h=552,\quad \Delta=331,\\
 k=668,\qquad \mu=249,\qquad R=\frac{551}{2},\qquad M=248.
\end{gathered}
\]
It has
\[
 R-\mu=\frac{53}{2}=26.5
 <\sqrt{999}-2+\frac1{500}.
\]
Nevertheless it obeys the uniform input, the degree-ratio bridge,
the path bound, and the maximum-degree broom bound. It also obeys the
weighted leaf-core upper bound obtained by replacing the nonleaf
core by a spanning tree:
\[
 \mu\le\min\left\{
 \frac{1001}{3},\
 \frac{4976587}{19980},\
 \frac{24886043}{99900}
 \right\}.
\]
The three right sides are respectively greater than $333$, $249.078$,
and $249.109$. The input requires only $M\ge247.753$, and the bridge
has right side greater than $472.818$, while $n-2R=449$.

Even the following two refinements do not exclude the point:
\[
\begin{split}
 n-2R\le {}&
 \left(\sqrt{\Delta/2}-1\right)(R-M)\\
 &+\left(1-\frac1{\sqrt2}\right)
          \left(1-\frac1{\sqrt\Delta}\right)L,\\
 R\ge {}& cI+
 L\left(\frac1{\sqrt\Delta}-\frac c\Delta\right),
 \qquad c=\frac{\sqrt{2\Delta}}{\Delta+2}.
\end{split}
\]
For the first, apply the minimum-degree-two bridge on nonleaf edges
and retain the correction on each leaf edge. For a leaf whose
neighbor has degree $d$, that correction is
\[
 (\sqrt d-\sqrt{\Delta/2})(d^{-1/2}-d^{-1}).
\]
Writing $x=d^{-1/2}$ makes this a convex quadratic on
$[1/\sqrt\Delta,1/\sqrt2]$, whose larger endpoint value occurs at
$d=\Delta$. This proves the displayed uniform leaf correction.
For the second inequality, a nonleaf edge has
\[
 \frac1{\sqrt{ab}}\ge c\left(\frac1a+\frac1b\right),
 \qquad 2\le a,b\le\Delta.
\]
If $P=\sum_{\ell v:\,d_\ell=1}1/d_v$, the core edges contribute at
least $c(I-P)$. The remaining term for each leaf is
$d^{-1/2}-c/d$, decreasing for $d\ge2$, and is at least its value
at $\Delta$.
At the proposed scalar point, the strengthened bridge has right side
greater than $450.005$, and the second bound requires only $R>67.192$.

\textbf{This point is not a constructed graph and is not a counterexample
to \textup{(CH)}.} It proves only that these explicitly listed aggregate
inequalities are jointly insufficient: they admit numerical data below
the target. A completion must retain additional graph structure,
particularly the interaction between leaf attachments and the cyclic
core. No logical recombination of these listed inequalities alone can exclude
this point; a stronger valid inequality or additional graph structure
is needed.

\subsection{The tree potential changes in the presence of cycles}

The tree degree-potential inequality cannot be applied unchanged to
graphs with cycles. Write $c=|E(G)|-n+1$ for the cyclomatic number and
retain $h=n-L-1$. The same edge potential now gives
\[
 R(G)\ge\sum_{d_v\ge2}\sqrt{d_v}
            -\left(\sqrt2-\frac12\right)(h+c),
 \qquad \sum_{d_v\ge2}(d_v-2)=L-2+2c.
\]
Consequently its chord bound becomes
\[
\begin{split}
 R(G)\ge{}&\frac h2+\sqrt\Delta+
                  \frac{L-\Delta}{\sqrt\Delta+\sqrt2}\\
 &+c\left\{\frac2{\sqrt\Delta+\sqrt2}
                         -\left(\sqrt2-\frac12\right)\right\}.
\end{split}
\]
The coefficient of $c$ is negative for $\Delta\ge3$. Any extension using
this potential therefore needs a corresponding gain in distance or a
stronger Randi\'c estimate; simply repeating the tree proof is insufficient.

\subsection{An exact formulation of the remaining core problem}

Delete the leaves of a mixed graph and let $H$ be the connected graph
induced by its nonleaves. For $i\in V(H)$ put
$a_i=d_H(i)$ and $w_i=1+\ell_i$, where $\ell_i$ is the number of
leaves formerly attached at $i$. Then
\[
 d_i=a_i+w_i-1\ge2,\qquad
 n=\sum_iw_i,\qquad L=\sum_i(w_i-1),
\]
and $H$ contains at least one cycle. Exact identities give
\[
 \begin{split}
 R(G)&=\sum_i\frac{w_i-1}{\sqrt{d_i}}
        +\sum_{ij\in E(H)}\frac1{\sqrt{d_id_j}},\\
 \mu(G)&=\frac n{n-1}\sum_{i,j}\frac{w_iw_j}{n^2}d_H(i,j)
                  +\frac{2L}{n}.
 \end{split}
\]
The second identity follows by separating the contribution of each leaf
edge; each contributes $n-1$ to the unordered distance sum.
Thus the unrestricted mixed case is exactly the assertion
\[
 \sum_i\frac{w_i-1}{\sqrt{d_i}}
 +\sum_{ij\in E(H)}\frac1{\sqrt{d_id_j}}
 -\frac n{n-1}\sum_{i,j}\frac{w_iw_j}{n^2}d_H(i,j)
 \ge\sqrt{n-1}-2+\frac{2(L+1)}n
\]
for such cyclic cores and integer weights. No proof of this full
weighted cyclic-core assertion is supplied here.

\section{Certificates, reproducibility, and provenance}

The proofs separate an infinite domain from a finite domain before
performing any computation. The finite checks use only integer or
rational arithmetic with outward square-root enclosures.
Floating point is used only for readable decimal displays.

\begin{center}\small
\begin{tabularx}{\linewidth}{@{}>{\raggedright\arraybackslash}p{.34\linewidth}
 r X@{}}
\toprule
Certificate & Cases & What is certified\\
\midrule
All trees, $n\ge1500$ & $10\,984$ &
Rational rectangles cover the entire reduced continuous domain\\
All trees, $16\le n\le1499$ & $552\,444\,094$ &
All feasible integer order, leaf-count and maximum-degree triples\\
Small-tree scalar reduction & $239$ &
All scalar cases through order $15$, leaving exactly seven profiles\\
Small-tree exceptional profiles & $4\,152$ &
Every subdivision allocation in the seven classified families\\
Leafless graphs, $3\le n\le255$ & $32\,131$ &
All integer order/maximum-degree pairs\\
General cutoff $3\,000\,000$ & $60$ &
All middle maximum-degree ratios via rational intervals\\
\bottomrule
\end{tabularx}
\end{center}

Stars and paths are handled explicitly. Leafless orders at least $256$
are handled analytically. The large-order general graph argument handles
the two outer maximum-degree ranges analytically as well.
Thus none of these results extrapolates from a finite sample.

The C++ finite-tree certificate uses signed 128-bit integers for every
proof comparison. Its arguments lie in the displayed bounded range;
the intermediate integer products, including cross-products used to
compare margins, have magnitude below $10^{31}$, well inside the signed
128-bit range. Square roots and inverse square roots are enclosed at
scale $10^{12}$ using integer binary search and division. The Python
certificates use unbounded integers and \texttt{fractions.Fraction}.

The tree estimates and code were independently checked through their
algebraic reductions and interval directions. The leafless broom
improvement and the three-million-order argument were also separately
audited. This is a reproducible research proof, including finite
certificates; it is not a claim of proof-assistant certification or
external referee approval.

For comparison, Cygan, Pilipczuk and \v{S}krekovski proved the earlier
tree estimate $R(T)-\mu(T)\ge\max\{0,\sqrt L-2\}$ \cite{cps}.
We do not infer historical priority for the present all-tree additive
bound from the searches conducted here. The frozen input used for the
other two theorems is VMRC commit
\texttt{52c1e4ae1f82bd3fdc8b809b9a94e2e4c6ef4cb3}.

\clearpage
\begin{thebibliography}{9}
\bibitem{cap}
G. Caporossi and P. Hansen,
Variable neighborhood search for extremal graphs 1. The AutoGraphiX system,
\emph{Discrete Mathematics} \textbf{212} (2000), 29--44.
\href{https://doi.org/10.1016/S0012-365X(99)00206-X}
{doi:10.1016/S0012-365X(99)00206-X}.

\bibitem{favaron}
O. Favaron, M. Mah\'eo and J.-F. Sacl\'e,
The Randi\'c index and other Graffiti parameters of graphs,
\emph{MATCH Communications in Mathematical and in Computer Chemistry}
\textbf{47} (2003), 7--23. Conjecture 4.8.
\href{https://match.pmf.kg.ac.rs/electronic_versions/Match47/match47_7-23.pdf}
{Journal PDF}.

\bibitem{vmrc227}
Verifiable Mathematical Reasoning Corpus,
\emph{A proof audit of Fajtlowicz's average-distance conjecture},
manuscript [227], 21 November 2026.
\href{https://github.com/QEEDDD/Verifiable-Mathematical-Reasoning-Corpus/blob/52c1e4ae1f82bd3fdc8b809b9a94e2e4c6ef4cb3/\%5B227\%5D\%20Fajtlowicz_Proof_Audit.tex}
{Frozen TeX source}. The imported assertion is its uniform specialization
of the weighted-distance theorem.

\bibitem{cps}
M. Cygan, M. Pilipczuk and R. \v{S}krekovski,
\emph{Relation between Randi\'c index and average distance of trees},
IMFM Preprint Series \textbf{48} (2010), no. 1130.
\href{https://users.fmf.uni-lj.si/skreko/Old/MyPapers/01130.pdf}
{Author-hosted primary PDF}.
\end{thebibliography}

\clearpage
\appendix
\section{Executable proof certificates}
\label{app:certificates}

The following sources are part of the proof, and are reproduced in
full so that this TeX file has no external source-code dependencies.
The accompanying verification archive supplies the same executable
files and records the executed results. The appendix listings need
not be executed to compile this document.




\subsection{All trees: the infinite-order reduction}
\texttt{ch\_tree\_infinite\_certificate.py}
\begin{lstlisting}[language=Python]
#!/usr/bin/env python3
"""Exact compact-domain certificate for the all-tree CH argument, n >= 1500.

All proof comparisons use fractions and integer square roots only.
"""
from fractions import Fraction as F
from math import isqrt
from pathlib import Path
import json


def sqrt_interval(x, scale=10**6):
    assert x >= 0
    a = isqrt(x.numerator * scale * scale // x.denominator)
    lo = F(a, scale)
    hi = lo if lo * lo == x else F(a + 1, scale)
    assert lo * lo <= x <= hi * hi
    return lo, hi


def run():
    n0 = 1500
    eps = F(37, 1000)
    strength = F(3869, 100)
    assert eps * eps > F(2, n0 - 1)
    assert strength * strength < F((n0 - 1)**3, n0**2)
    assert F(21, 10) * (1 - eps) > 2
    assert strength / 24 > 1
    smallest = None
    witness = None
    count = 0
    for iq in range(110):
        ql, qu = F(100 + iq, 100), F(101 + iq, 100)
        for iz in range(100):
            zl, zu = F(iz, 200), F(iz + 1, 200)
            if ql * zl >= 1:
                continue
            rl = sqrt_interval(max(F(0), 1 - qu * zu))[0]
            ru = sqrt_interval(1 - ql * zl)[1]
            a_lower = (ql - 1) / (ru + eps) - qu / (1 + rl)
            b_lower = zl * zl / 6 + max(F(0), 1 - (2 * qu + 1) * zu)**2 / 2
            margin = a_lower + strength * b_lower
            assert margin > 0, (ql, qu, zl, zu, margin)
            count += 1
            if smallest is None or margin < smallest:
                smallest = margin
                witness = [ql, qu, zl, zu]
    result = {
        'status': 'passed',
        'scope': 'All connected trees of every order n >= 1500, after the displayed reductions',
        'rational_rectangle_count': count,
        'epsilon_upper': str(eps),
        'strength_lower': str(strength),
        'minimum_certified_margin': str(smallest),
        'minimum_margin_decimal_for_display_only': float(smallest),
        'witness_rectangle_qlo_qhi_zlo_zhi': [str(x) for x in witness],
        'outer_q_margin': str(F(21, 10) * (1 - eps) - 2),
        'outer_z_margin': str(strength / 24 - 1),
        'sqrt_enclosure_scale': 10**6,
    }
    target = Path(__file__).with_suffix('.json')
    target.write_text(json.dumps(result, indent=2) + '\n')
    print(json.dumps(result, indent=2))


if __name__ == '__main__':
    run()

\end{lstlisting}



\subsection{All trees: orders 16 through 1499}
\texttt{ch\_tree\_finite\_certificate.cpp}
\begin{lstlisting}[language=C++]
#include <algorithm>
#include <cstdint>
#include <cstdlib>
#include <iomanip>
#include <iostream>
#include <limits>
#include <vector>

// All proof decisions use signed __int128 integers. Floating point is used
// solely to print the minimum certified margin after the exact comparisons.
using i128 = __int128_t;
using i64 = std::int64_t;
constexpr i64 S = 1000000000000LL;
constexpr int ROOT_LIMIT = 4000;
std::vector<i64> rl(ROOT_LIMIT + 1), ru(ROOT_LIMIT + 1);
std::vector<i64> il(ROOT_LIMIT + 1), iu(ROOT_LIMIT + 1);

void initialize_intervals() {
    for (int j = 1; j <= ROOT_LIMIT; ++j) {
        const i128 target = i128(j) * S * S;
        i64 lo = 0, hi = 100 * S;
        while (hi - lo > 1) {
            i64 mid = lo + (hi - lo) / 2;
            if (i128(mid) * mid <= target) lo = mid;
            else hi = mid;
        }
        rl[j] = lo;
        ru[j] = lo + 1;
        il[j] = i128(S) * S / ru[j];
        iu[j] = (i128(S) * S + rl[j] - 1) / rl[j];
        if (!(i128(rl[j])*rl[j] <= target &&
              target < i128(ru[j])*ru[j])) std::abort();
    }
}

int main(int argc, char** argv) {
    int nmin = argc > 1 ? std::atoi(argv[1]) : 16;
    int nmax = argc > 2 ? std::atoi(argv[2]) : 1499;
    if (nmin < 4 || nmax >= 2000 || nmin > nmax) return 2;
    initialize_intervals();
    std::uint64_t checked = 0, failed = 0;
    i128 min_margin = 0;
    i64 min_den = 1;
    int min_n = 0, min_h = 0, min_D = 0;
    for (i64 n = nmin; n <= nmax; ++n) {
        const i64 N = n - 1;
        const i64 den = 6*n*N;
        for (i64 h = 1; h <= n-3; ++h) {
            const i64 I = h+1, L = N-h;
            // An internal degree sequence is 2 plus nonnegative excesses,
            // totaling L-2, each at most Delta-2.
            const i64 Dmin = std::max<i64>(3, 2+(L-2+I-1)/I);
            const i64 Rleaf = h*S/2 + rl[L] + il[2] - S/2
                              - iu[L] + il[2*L];
            const i64 nroot = ru[N];
            for (i64 D = Dmin; D <= L; ++D) {
                const i64 q = (L-2)/(D-2), rem = (L-2)%(D-2);
                const i64 e = rem != 0;
                if (q+e > I) continue;
                const i64 Rdegree = h*S/2 + q*rl[D]
                    + (e ? rl[rem+2] : 0) + (1-q-e)*ru[2];
                const i64 Rlower = std::max(Rdegree, Rleaf);
                const i64 k = N-D;
                const i64 a = std::max<i64>(
                    D-h+1, 1+(L+I-1)/I);
                const i64 delta = std::max<i64>(
                    0, 2*a-(n-h+1));
                const i64 core_num =
                    3*h*n*n-h*(h*h-1)-3*h*delta*delta+12*L*N;
                const i64 broom_num =
                    12*N*N+2*k*(k+1)*(3*n-2*k-7);
                const i64 path_num = 2*n*N*(n+1);
                const i64 mu_num = std::min({
                    core_num,broom_num,path_num});
                // Rlower/S - mu_num/den
                //   - [sqrt(n-1)-2+2/n] > 0.
                const i128 margin =
                    i128(Rlower-nroot+2*S)*den
                    -i128(S)*(mu_num+12*N);
                ++checked;
                if (min_n == 0 ||
                    margin*min_den < min_margin*den) {
                    min_margin = margin;
                    min_den = den;
                    min_n = n; min_h = h; min_D = D;
                }
                if (margin <= 0) {
                    ++failed;
                    if (failed <= 20) {
                        std::cout << "FAIL n=" << n << " h=" << h
                                  << " L=" << L << " Delta=" << D
                                  << " a=" << a
                                  << " lower_slack="
                                  << std::setprecision(15)
                                  << (long double)margin/(S*(long double)den)
                                  << '\n';
                    }
                }
            }
        }
    }
    std::cout << "Range " << nmin << ".." << nmax
              << "; checked=" << checked
              << "; failures=" << failed << '\n';
    std::cout << "Minimum certified slack="
              << std::setprecision(18)
              << (long double)min_margin/(S*(long double)min_den)
              << "; n=" << min_n << "; h=" << min_h
              << "; L=" << (min_n-1-min_h)
              << "; Delta=" << min_D << '\n';
    return failed ? 1 : 0;
}

\end{lstlisting}



\subsection{All trees: the remaining small profiles}
\texttt{ch\_tree\_small\_certificate.py}
\begin{lstlisting}[language=Python]
"""Exact certificate for seven exceptional small tree parameter profiles.

All vertices of degree at least three are the two hubs, with degrees
(4,3) or (5,3). The intervening vertices of degree two are distributed
among the hub connector and the pendant arms. Unordered-arm symmetry
is intentionally not quotiented: coverage is simpler with all ordered
weak compositions.
"""
from fractions import Fraction as Q
from math import isqrt
from functools import cache

S = 10**12
PARAMETERS = (
    (11,5,4), (12,6,5), (12,5,4), (13,6,5),
    (13,5,4), (14,6,5), (15,6,5),
)

@cache
def root_upper(k):
    return Q(isqrt(k*S*S)+1,S)

@cache
def root_lower(k):
    return Q(isqrt(k*S*S),S)

@cache
def inverse_root_lower(k):
    return 1/root_upper(k)

@cache
def inverse_root_upper(k):
    return 1/root_lower(k)

scalar_count = 0
scalar_exceptions = []
for n in range(4,16):
    N = n-1
    for h in range(1,n-2):
        I,L = h+1,N-h
        for D in range(max(3,2+(L-2+I-1)//I),L+1):
            q,rem = divmod(L-2,D-2)
            e = int(rem != 0)
            if q+e > I:
                continue
            Ra = (Q(h,2)+q*root_lower(D)
                  +(root_lower(rem+2) if e else 0)
                  +(1-q-e)*root_upper(2))
            Rb = (Q(h,2)+root_lower(L)+inverse_root_lower(2)
                  -Q(1,2)-inverse_root_upper(L)
                  +inverse_root_lower(2*L))
            k = N-D
            a = max(D-h+1,1+(L+I-1)//I)
            delta = max(0,2*a-(n-h+1))
            den = 6*n*N
            core = Q(3*h*n*n-h*(h*h-1)-3*h*delta*delta+12*L*N,den)
            broom = Q(12*N*N+2*k*(k+1)*(3*n-2*k-7),den)
            path = Q(n+1,3)
            margin = (max(Ra,Rb)-min(core,broom,path)
                      -root_upper(N)+2-Q(2,n))
            scalar_count += 1
            if margin <= 0:
                scalar_exceptions.append((n,L,D))
assert scalar_count == 239
assert sorted(scalar_exceptions) == sorted(PARAMETERS)
print("Small scalar cases:",scalar_count,
      "| exactly seven remaining profiles:",scalar_exceptions)

def compositions(total, slots):
    if slots == 1:
        yield (total,)
    else:
        for first in range(total+1):
            for rest in compositions(total-first,slots-1):
                yield (first,)+rest

def build_tree(p, extras):
    edges = []
    next_id = 2
    # Hubs are 0 and 1. Connector has 1+extras[0] edges.
    current = 0
    for _ in range(extras[0]):
        edges.append((current,next_id))
        current = next_id
        next_id += 1
    edges.append((current,1))
    # Hub 0 has p-1 pendant arms; hub 1 has 2 pendant arms.
    hubs = [0]*(p-1)+[1]*2
    assert len(hubs) == len(extras)-1
    for hub, extra in zip(hubs,extras[1:]):
        current = hub
        for _ in range(extra+1):
            edges.append((current,next_id))
            current = next_id
            next_id += 1
    return next_id, edges

def certified_slack(n, edges):
    adj = [[] for _ in range(n)]
    for u,v in edges:
        adj[u].append(v)
        adj[v].append(u)
    degree = list(map(len,adj))
    Rlower = sum(
        (inverse_root_lower(degree[u]*degree[v]) for u,v in edges),
        Q(0),
    )
    parent = [-1]*n
    parent[0] = 0
    order = [0]
    for v in order:
        for u in adj[v]:
            if u != parent[v]:
                assert parent[u] == -1
                parent[u] = v
                order.append(u)
    assert len(order) == n
    size = [1]*n
    W = 0
    for v in reversed(order[1:]):
        W += size[v]*(n-size[v])
        size[parent[v]] += size[v]
    return (Rlower-Q(2*W,n*(n-1))-root_upper(n-1)
            +2-Q(2,n))

checked = 0
minimum = None
witness = None
counts = []
for n,L,p in PARAMETERS:
    local_count = 0
    for extra in compositions(n-L-2,L+1):
        count,edges = build_tree(p,extra)
        assert count == n and len(edges) == n-1
        lower = certified_slack(n,edges)
        assert lower > 0, (n,L,p,extra,lower)
        checked += 1
        local_count += 1
        if minimum is None or lower < minimum:
            minimum = lower
            witness = (n,L,p,extra)
    counts.append(local_count)
assert counts == [126,210,252,462,462,924,1716]
assert checked == 4152
print("Exact small tree profiles passed:", checked)
print("Counts by parameter triple:", counts)
print("Minimum certified slack:", float(minimum))
print("Witness:", witness)

\end{lstlisting}



\subsection{All graphs without leaves: complete certificate}
\texttt{ch\_min\_degree\_two\_certificate.py}
\begin{lstlisting}[language=Python]
#!/usr/bin/env python3
"""Exact certificate for the minimum-degree-two CH reduction.

The theorem uses the corpus uniform input CH-U, the degree-ratio bridge,
and the separately proved leafless improvement of the broom distance bound.
Only parameter inequalities are checked: no claim of graph enumeration.
"""

from fractions import Fraction as F
from math import isqrt
from pathlib import Path
import json

SCALE = 1_000_000


def sqrt_upper(x):
    x = F(x)
    scaled = x * SCALE * SCALE
    r = isqrt(scaled.numerator // scaled.denominator)
    if F(r * r, SCALE * SCALE) < x:
        r += 1
    answer = F(r, SCALE)
    assert answer * answer >= x
    return answer


def broom_bound(n, d):
    k = n - 1 - d
    return 2 - F(2, n) + F(k * (k + 1) * (3 * n - 2 * k - 7), 3 * n * (n - 1))


def main():
    a, c0_bound = F(71, 100), F(231, 1000)
    assert a * a > F(1, 2)
    assert (1 - 2 * c0_bound) ** 2 < 1 - a
    critical_value_upper = F(2, 3) + c0_bound + c0_bound**2 - F(8, 3) * c0_bound**3
    assert critical_value_upper < F(23, 25)
    assert F(2, 3) + a / 4 < F(23, 25)
    assert 5 * a / 4 < F(23, 25)
    large_order_upper = F(23, 25) + F(4, 255) + F(1, 16)
    assert large_order_upper < 1

    per_order = []
    case_count = 0
    global_minimum = None
    for n in range(3, 256):
        s = sqrt_upper(n - 1)
        gamma_upper = s - 2 + F(2, n)
        worst = None
        for d in range(2, n):
            k = n - 1 - d
            B = broom_bound(n, d) - F(2 * k, n * (n - 1))
            r = sqrt_upper(F(d, 2))
            rhs_upper = 2 * (B + gamma_upper) + (r - 1) * (s - 1 + B / n)
            margin = n - rhs_upper
            assert margin > 0, (n, d, str(margin))
            case_count += 1
            record = {
                "n": n,
                "maximum_degree": d,
                "k": k,
                "mean_distance_upper_exact": str(B),
                "sqrt_n_minus_1_upper": str(s),
                "sqrt_degree_over_2_upper": str(r),
                "strict_margin_exact": str(margin),
                "strict_margin_decimal_for_display": float(margin),
            }
            if worst is None or margin < F(worst["strict_margin_exact"]):
                worst = record
            if global_minimum is None or margin < F(global_minimum["strict_margin_exact"]):
                global_minimum = record
        per_order.append(worst)
    assert case_count == 32131
    payload = {
        "status": "passed",
        "theorem_scope": "Every finite connected simple graph with minimum degree at least two, conditional on CH-U and the proved leafless broom lemma",
        "uniform_input": "M >= ((n-1)*mu - (n-2))/n",
        "degree_ratio_bridge": "n - 2R <= (sqrt(Delta/2)-1)*(R-M)",
        "leafless_broom_improvement": "mu <= B(n,Delta) - k/binom(n,2), k=n-1-Delta",
        "analytic_all_large_orders": {
            "n_minimum": 256,
            "bound_for_2b_plus_sqrt_c_over_2": "23/25",
            "critical_value_upper_exact": str(critical_value_upper),
            "final_normalized_upper_exact": str(large_order_upper),
            "final_normalized_margin_exact": str(1 - large_order_upper),
        },
        "finite_parameter_check": {
            "order_minimum": 3,
            "order_maximum": 255,
            "case_count": case_count,
            "sqrt_upper_denominator": SCALE,
            "global_minimum": global_minimum,
            "worst_case_per_order": per_order,
        },
        "qualification": "The finite check is exhaustive in the stated (n,Delta) parameter range. The analytic argument handles all n>=256. Graph reduction and input validation are mathematical dependencies, not established merely by running this script.",
    }
    out = Path(__file__).resolve().parent / "ch_min_degree_two_certificate.json"
    out.write_text(json.dumps(payload, indent=2) + "\n", encoding="utf-8")
    print(json.dumps({
        "status": "passed",
        "all_large_orders_from": 256,
        "finite_parameter_cases": case_count,
        "global_minimum": global_minimum,
        "large_order_normalized_margin": str(1-large_order_upper),
    }, indent=2))


if __name__ == "__main__":
    main()

\end{lstlisting}



\subsection{The three-million-order cutoff}
\texttt{ch\_heavy\_cutoff\_certificate.py}
\begin{lstlisting}[language=Python]
#!/usr/bin/env python3
"""Exact rational interval certificate for the conditional CH cutoff 3,000,000.

This certifies inequalities in a separately derived mathematical reduction.
It does not enumerate graphs and is not an unconditional proof of the graph
theorem: the reduction uses the corpus uniform weighted-distance input (CH-U).
"""

from fractions import Fraction as F
from math import isqrt
from pathlib import Path
import json

N = 3_000_000
EPS = F(1, N)
U = F(1, 144)       # n^(-1/3) < U for n >= N
V = F(1, 1732)      # n^(-1/2) < V for n >= N
SQRT7_UPPER = F(2646, 1000)
A_UPPER = F(9000003, 200000)  # 45.000015
SCALE = 100_000


def sqrt_upper(x):
    scaled = x * SCALE * SCALE
    numerator = isqrt(scaled.numerator // scaled.denominator)
    if F(numerator * numerator, SCALE * SCALE) < x:
        numerator += 1
    answer = F(numerator, SCALE)
    assert answer * answer >= x
    return answer


def positive(x):
    return max(F(0), x)


def f(z, c):
    return z**3 + 3 * z * positive(2 * c + z - 1) ** 2


def phi(z, c):
    return z / 2 - f(z, c) / 6


def E(c):
    return (
        (3 * SQRT7_UPPER + A_UPPER * (1 - c) / 7) * U
        + 3 * V
        + A_UPPER * EPS
        + EPS**2
    )


def main():
    assert 144**3 < N
    assert 1732**2 < N - 1
    assert SQRT7_UPPER**2 > 7
    assert A_UPPER > 45 + F(42, N - 1)
    records = []
    for j in range(10, 70):
        left, right = F(j, 100), F(j + 1, 100)
        bound = E(left)
        lo, hi = 0, int((1 - left) * SCALE)
        assert f(F(hi, SCALE), left) > bound
        while lo + 1 < hi:
            middle = (lo + hi) // 2
            if f(F(middle, SCALE), left) > bound:
                hi = middle
            else:
                lo = middle
        H = F(hi, SCALE)
        assert 0 < H <= 1 - left
        assert f(H, left) > bound
        mean_upper = (
            phi(H, left) + EPS**2 / 6 + 7 * EPS + (1 - left) * U
        ) / (1 - EPS)
        mean_lower = (1 - sqrt_upper(right)) / 2 - V / 2
        margin = mean_lower - mean_upper
        assert margin > 0, (j, float(margin))
        records.append(
            {
                "c_lower": str(left),
                "c_upper": str(right),
                "H_upper": str(H),
                "sqrt_c_upper_bound": str(sqrt_upper(right)),
                "fH_minus_E_exact": str(f(H, left) - bound),
                "mean_upper_exact": str(mean_upper),
                "mean_lower_exact": str(mean_lower),
                "strict_margin_exact": str(margin),
                "mean_upper_decimal_for_display": float(mean_upper),
                "mean_lower_decimal_for_display": float(mean_lower),
                "margin_decimal_for_display": float(margin),
            }
        )

    low_degree_mean_upper = (1 + EPS) / 3
    low_degree_mean_lower = (1 - sqrt_upper(F(1, 10))) / 2 - V / 2
    assert low_degree_mean_lower > low_degree_mean_upper
    high_degree_forced_rhs = 2 * V + (2 + V) * (F(6, 25) + 2 * EPS)
    high_degree_forced_lhs = F(1, 2) - V / 2
    assert high_degree_forced_lhs > high_degree_forced_rhs
    assert F(6, 25) + 2 * EPS < 1  # used to bound the final fraction

    output = {
        "status": "passed",
        "cutoff": N,
        "conditional_input": "(CH-U): M >= ((n-1)*mu-(n-2))/n",
        "variables": "c=Delta/n, z=h/n; selected high centers include a max-degree vertex; K=ceil(7*n^(1/3))",
        "bounds": {
            "n_inverse": str(EPS),
            "n_inverse_cube_root": str(U),
            "n_inverse_square_root": str(V),
            "sqrt_7": str(SQRT7_UPPER),
            "45_plus_42_over_n_minus_1": str(A_UPPER),
        },
        "interval_count": len(records),
        "minimum_margin_decimal_for_display": min(
            record["margin_decimal_for_display"] for record in records
        ),
        "low_degree_margin_exact": str(low_degree_mean_lower - low_degree_mean_upper),
        "high_degree_margin_exact": str(high_degree_forced_lhs - high_degree_forced_rhs),
        "records": records,
        "scope": "Rational verification of parameter inequalities after the stated graph reduction; not graph enumeration or unconditional validation of CH-U",
    }
    destination = Path(__file__).resolve().parent
    (destination / "ch_heavy_cutoff_certificate.json").write_text(
        json.dumps(output, indent=2) + "\n", encoding="utf-8"
    )
    table = [
        "% Generated by ch_heavy_cutoff_certificate.py; displayed decimals are rounded.",
        "\\begin{longtable}{rrrrr}",
        "\\toprule",
        "$c_-$ & $c_+$ & $H$ & upper bound for $\\mu/n$ & lower bound for $\\mu/n$\\\\",
        "\\midrule\\endhead",
    ]
    for record in records:
        left = F(record["c_lower"])
        right = F(record["c_upper"])
        H = F(record["H_upper"])
        # Display upper bounds rounded upward, lower bounds rounded downward.
        upper = F(record["mean_upper_exact"])
        lower = F(record["mean_lower_exact"])
        shown_upper = F((upper.numerator * 10**6 + upper.denominator - 1) // upper.denominator, 10**6)
        shown_lower = F(lower.numerator * 10**6 // lower.denominator, 10**6)
        assert shown_upper < shown_lower
        table.append(
            f"{float(left):.2f} & {float(right):.2f} & {float(H):.5f} & "
            f"{float(shown_upper):.6f} & {float(shown_lower):.6f}\\\\"
        )
    table += ["\\bottomrule", "\\end{longtable}", ""]
    (destination / "ch_heavy_cutoff_table.tex").write_text("\n".join(table), encoding="utf-8")
    print(json.dumps({
        "status": output["status"],
        "cutoff": N,
        "interval_count": len(records),
        "minimum_margin": output["minimum_margin_decimal_for_display"],
        "low_degree_margin": float(low_degree_mean_lower - low_degree_mean_upper),
        "high_degree_margin": float(high_degree_forced_lhs - high_degree_forced_rhs),
    }, indent=2))


if __name__ == "__main__":
    main()

\end{lstlisting}



\subsection{The scalar relaxation witness}
\texttt{ch\_mixed\_scalar\_witness.py}
\begin{lstlisting}[language=Python]
#!/usr/bin/env python3
"""Exact scalar feasibility check; this tuple does not assert a graph exists."""
from fractions import Fraction as F
from math import isqrt

S = 10**12

def sqrt_interval(x):
    x = F(x)
    q = isqrt((x.numerator*S*S)//x.denominator)
    lo, hi = F(q,S), F(q+1,S)
    assert lo*lo <= x < hi*hi
    return lo, hi

n, L, I, h, Delta, k = 1000, 447, 553, 552, 331, 668
N = n-1
mu, R, M = F(249), F(551,2), F(248)
X = R-M
rootD_lo, rootD_hi = sqrt_interval(Delta)
root2_lo, root2_hi = sqrt_interval(2)
ratio_lo, ratio_hi = sqrt_interval(F(Delta,2))
rootN_lo, rootN_hi = sqrt_interval(N)
root2D_lo, root2D_hi = sqrt_interval(2*Delta)

a = max(Delta-h+1, 1+(L+I-1)//I)
imbalance = max(0, 2*a-(n-h+1))
mu_broom = F(2)-F(2,n)+F(k*(k+1)*(3*n-2*k-7),3*n*N)
mu_heavy = F(3*h*n*n-h*(h*h-1)-3*h*imbalance**2+12*L*N,6*n*N)
M_uniform = F(N*mu-(n-2),n)
M_deleted = F(n*mu-F(k*(k+1),N),N)-1

assert mu <= F(n+1,3)
assert mu <= mu_broom and mu <= mu_heavy
assert M >= M_uniform and M >= M_deleted
assert rootN_hi <= R <= F(n,2) and M <= R
assert 0 <= X <= R-mu+1+F(mu-2,n)
assert n-2*R <= (rootD_lo-1)*X
leaf_correction_lo = L*(1-1/root2_lo)*(1-1/rootD_lo)
refined_rhs_lo = (ratio_lo-1)*X+leaf_correction_lo
assert n-2*R <= refined_rhs_lo
mixed_R_upper = F(root2D_hi,Delta+2)*(I-F(L,Delta))+F(L,rootD_lo)
assert R >= mixed_R_upper
diameter = 300
rad_lo, rad_hi = sqrt_interval(n-diameter+1)
assert mu <= diameter <= n-Delta+1
assert M >= F(diameter,2)
assert R >= F(diameter,2)+rad_hi-1
assert R-mu < rootN_lo-2+F(2,n)

print('PASS: an abstract scalar point below the target satisfies every listed bound.')
print('This is not an actual graph or a graph counterexample.')
print('Gap:', R-mu, '; target lower bound:', float(rootN_lo-2+F(2,n)))
print('Refined bridge lower bound:', float(refined_rhs_lo), '; left side:', n-2*R)
print('Broom mean upper:', mu_broom, '; heavy-core mean upper:', mu_heavy)

\end{lstlisting}


\end{document}
